\documentclass[10pt,letterpaper]{article}
\usepackage[T1]{fontenc}
\usepackage{lmodern}
\usepackage{microtype}
\usepackage[margin=0.83in,headheight=14pt]{geometry}
\usepackage{amsmath,amssymb,amsthm,mathtools}
\usepackage{booktabs,longtable,array}
\usepackage{xcolor}
\usepackage{enumitem}
\usepackage{fancyhdr}
\usepackage{hyperref}
\definecolor{inkblue}{HTML}{173B57}
\hypersetup{colorlinks=true,linkcolor=inkblue,urlcolor=inkblue,citecolor=inkblue,
 pdftitle={Pell Centers and Fused Bounds for Diophantine Consistency},
 pdfauthor={Research note prepared for Vladimir Reshetnikov, with AI assistance}}
\pagestyle{fancy}\fancyhf{}
\fancyhead[L]{\small\textsc{ProveIt / Diophantine consistency}}
\fancyhead[R]{\small 25 October 2026}
\fancyfoot[C]{\thepage}
\setlength{\parskip}{4pt}
\setlength{\parindent}{0pt}
\setlist{itemsep=2pt,topsep=4pt,leftmargin=*}
\newtheorem{theorem}{Theorem}[section]
\newtheorem{lemma}[theorem]{Lemma}
\newtheorem{proposition}[theorem]{Proposition}
\newtheorem{corollary}[theorem]{Corollary}
\theoremstyle{remark}\newtheorem{remark}[theorem]{Remark}
\newcommand{\Z}{\mathbb Z}\newcommand{\N}{\mathbb N}
\newcommand{\Con}{\operatorname{Con}}
\newcommand{\ZFC}{\mathsf{ZFC}}
\newcommand{\ZF}{\mathsf{ZF}}
\newcommand{\vp}{\nu_2}
\newcommand{\pathtext}[1]{\nolinkurl{#1}}
\allowdisplaybreaks[2]

\begin{document}
\begin{center}
{\Large\bfseries\color{inkblue} Pell Centers and Fused Bounds}\\[5pt]
{\large Explicit Diophantine presentations of $\Con(\ZFC)$}\\[9pt]
Prepared for Vladimir Reshetnikov and the ProveIt research program\\
Research and artifact development with AI assistance\\[4pt]
25 October 2026
\end{center}

\begin{abstract}
We give three arithmetic reductions of explicit consistency equations in
the pinned ProveIt snapshot. A polynomial Pell center replaces a large
constant in the growth certificate and yields a quartic in 212 unrestricted
integer variables. Its small-coefficient version has 152 quadratic residual
squares, 21,556 collected monomials, and coefficient height 6,432; a sparse
version has 5,290 monomials and height 277,420. A certificate shared between
the Pell-index cap and the binomial-remainder bound removes three integer
variables. Combining it with a linear Pell-orientation numerator gives an
82-variable equation of exact degree 16,964 with 1,723 live binary arithmetic
operations; companions have 81 variables at degree 2,097,732 and 106 variables
at degree 32. Separately, an odd-index Pell residue eliminates a native
transport quotient, giving a 13-positive-variable equation of degree 174
with a 69-operation evaluator.

Complete new ordinary proofs, literal equations, finite algebraic certificates,
independent checkers, and reproducible construction programs accompany this
report. Composed with the explicitly identified inherited logical
calibrations, the proofs are formalizable in ZFC and establish that the
respective equations have a zero exactly when ZFC has a finite refutation.
Global minimality and a fresh end-to-end proof-assistant compilation remain
open.
\end{abstract}

\setcounter{tocdepth}{1}
\tableofcontents
\newpage

\section{Results and their exact scope}

Throughout, $\Con(\ZFC)$ means absence of a finite syntactic refutation in
the fixed effective ZFC presentation used by ProveIt. Effective translations
between the standard presentations give the usual consistency statement.
All polynomial coefficients are integers. A named register in a circuit is
an abbreviation for a polynomial, never an additional quantified variable.

\begin{theorem}[Consistency equations]
For each integer-domain polynomial $P$ in Table~\ref{tab:new}, with its
indicated number $n$ of inputs, the ordinary proofs in this report,
composed with the inherited calibration identified in
Section~\ref{sec:calibration}, establish
\[
 \ZFC\vdash
 \bigl(\neg\exists\boldsymbol x\in\Z^n\;P(\boldsymbol x)=0\bigr)
 \longleftrightarrow\Con(\ZFC).
\]
For the positive-domain polynomials, replace $\Z^n$ by
$\Z_{>0}^{13}$. Their explicit 39-integer lifts have the same
nonsolvability equivalence over $\Z^{39}$.
\end{theorem}

\begin{table}[ht]
\centering\small
\begin{tabular}{@{}llrrl@{}}
\toprule
Polynomial & Domain & Inputs & Exact degree & Further size data\\
\midrule
$P_{212}^{\rm small}$ & $\Z$ &212&4&152 squares; height 6,432\\
$P_{212}^{\rm sparse}$ & $\Z$ &212&4&5,290 monomials; height 277,420\\
$P_{106}$ & $\Z$ &106&32&47 squares; 1,830 operations\\
$P_{82}$ & $\Z$ &82&16,964&23 squares; 1,723 operations\\
$P_{81}$ & $\Z$ &81&2,097,732&22 squares; 1,711 operations\\
$C_{13}$ & $\Z_{>0}$ &13&174&69 operations\\
$C_{13}^{\rm oriented}$ & $\Z_{>0}$ &13&170&70 operations\\
$\widetilde C_{39}$ & $\Z$ &39&348&160 operations\\
$\widetilde C_{39}^{\rm oriented}$ & $\Z$ &39&340&161 operations\\
\bottomrule
\end{tabular}
\caption{Explicit equations supplied in this submission. Operations count
live binary $+,-,\times$ gates with integer leaves.}
\label{tab:new}
\end{table}

The principal quartic is literally
\begin{equation}\label{eq:mainquartic}
 P_{212}^{\rm small}(x_0,\ldots,x_{211})
 =\sum_{j=0}^{151}R_j(x_0,\ldots,x_{211})^2=0,
 \qquad (x_0,\ldots,x_{211})\in\Z^{212}.
\end{equation}
Every $R_j$ is printed in Appendix~\ref{app:literal}, and the JSON file
\pathtext{quartic212/reduced/equation.json} gives the same 152 literal
quadratic polynomials. The full collected quartic is also supplied. Thus
the definition contains no implicit universal polynomial, unspecified
search procedure, or variable exponent.

The high-degree direct equations are supplied as acyclic arithmetic graphs,
with independent input lists and fixed residual outputs. Their specification
is just as finite and literal, while avoiding expansion of very large
polynomials. The 13-positive-variable formula is short enough to print in
full in the main text.

\subsection{What improves, and what remains a tradeoff}

The inspected repository revision is
\pathtext{b42c796d1f5de9912a576bc2fd77b778ae7402bf}.
It includes the latest incoming archives, not only the current root display.
The comparisons below refer to actual domain and arithmetic data.

\begin{table}[ht]
\centering\small
\begin{tabular}{@{}lrrrr@{}}
\toprule
Matched direct family & Integer inputs & Degree & Operations & Squares\\
\midrule
Incoming projected source &109&32&1,895&48\\
New fused and oriented &106&32&1,830&47\\
Incoming projected source &85&16,964&1,788&24\\
New fused and oriented &82&16,964&1,723&23\\
Incoming projected source &84&2,097,732&1,775&23\\
New fused and oriented &81&2,097,732&1,711&22\\
\bottomrule
\end{tabular}
\caption{The three-variable reduction preserves the matched exact degree.}
\end{table}

The incoming growth quartic~\cite{growth} has 210 variables, 150 squares, 5,282 monomials,
and coefficient height $2^{40}$. Our polynomial-center replacement uses two
additional variables and two additional squares. After exact integral row
reduction, it has height 6,432. The sparse companion preserves a much smaller
support at the larger height 277,420. An existing balanced 214-variable
quartic has height 91,679, whereas recent 215-variable proposals reach
height 2,304. Those alternatives remain useful. The new quartic improves
the low-variable, small-coefficient tradeoff; it does not subsume every
existing point on the frontier.

The new positive carrier is compared with the exact incoming Carrier65~\cite{carrier}:
14 positive variables, degree 150, 65 operations. It removes one variable
at the cost of four operations and degree 174. Its alternative lowers the
degree to 170 with one further operation. Other repository equations have
fewer variables, different degrees, or fewer operations. In particular,
81 is not a global variable-count record, and the 39-integer translation
of the compact equation is a transparent domain conversion rather than
a claim of optimality.

The phrase ``simplest equation'' needs a chosen metric. Variables, degree,
expanded support, coefficient height, circuit size, coefficient-description
length, and proof length can move in different directions. This report
records those distinctions instead of assigning an unsupported universal
ordering. The new ideas are concrete improvements relative to the inspected
families. A claim of historical priority across all literature would require
a separate investigation.

\section{The arithmetic and logical dependency boundary}\label{sec:calibration}

The two direct families use the repository's numerical calculus C222.
Their semantic chain is
\[
\begin{gathered}
\text{integer zero}\quad\Longleftrightarrow\quad
\text{finite square-shell C222 refutation}\\
\Longleftrightarrow\quad\text{finite original C222 refutation}
\quad\Longleftrightarrow\quad\text{finite ZF refutation}\\
\Longleftrightarrow\quad\text{finite ZFC refutation}.
\end{gathered}
\]
The square-shell recoding is an explicit recursive bijection on numerical
syntax. It fixes truth 0 and falsity 1 and transports the guarded schema.
The next two logical endpoints are the identified repository results
\pathtext{KEndpoint.c222_iff_zf_refutation} and
\pathtext{ConstructibleChoice.zf_refutation_iff_zfc_refutation}.
Their source review and exact provenance are retained in the archive~\cite{calibration}.
The direct equations do not require a Turing-machine simulation or Jones'
universal polynomial. The separate compact family retains its native
coefficient/compiler ancestry, including the inherited machine-to-refutation
calibration.

On the arithmetic side, the inherited direct source proves finite
Horn-flow extraction, beta-CRT reconstruction, factorial-free binomial
selection, and signed Pell index classification. This submission changes
three specifically identified arithmetic interfaces and proves their new
soundness and completeness. It does not infer those unbounded facts from
sample computations.

The proof status has three distinct parts. First, exact finite polynomial
identities and source bindings are checked by executable verifiers. Second,
the new unbounded arithmetic lemmas are ordinary mathematical proofs,
reviewed independently in this work. Third, the C222/ZF/ZFC and native
calibrations are credited inherited results. The selected Lean source
evidence is not a full newly compiled transitive import closure, and this
release makes no new end-to-end Lean or Rocq claim.

All recursions and completeness constructions concern finite internally
coded syntax, natural numbers, finite rings, and finite lists. Their
induction, integer arithmetic, finite CRT, and Pell arguments are available
in ZFC. Consequently the intended biconditional concerns each model's
internal integers and finite derivations. It does not impose an external
``standard proof code'' restriction or assume omega-consistency. The
equivalence proves neither $\Con(\ZFC)$ nor its negation.

\section{A polynomial Pell center with small coefficients}

\subsection{The growth interface}

The square-shell constructors used by the direct source are
\begin{align*}
 I(u,v)&=2((u+v)^2+u),& M(u,v)&=I(u,v)+1,\\
 A(u,v)&=2((u+v)^2+3u+4v+3),& E(u,v)&=A(u,v)+1.
\end{align*}
The authenticated local axiom terms have depths
$4,3,3,3,3,5,4,3,19$. Every local leaf parameter is below the horizon $t$.
For $0\le u,v\le m$, the largest constructor satisfies
\[
 E(m,m)+1=8m^2+14m+8\le8(m+1)^2.
\]
Induction on depth therefore bounds every local formula label strictly by
\begin{equation}\label{eq:labelbound}
 L=8^{d-1}t^d,\qquad d=2^{19}.
\end{equation}
Indeed the recurrence $H_0=t$, $H_{j+1}=8H_j^2$ has solution
$H_j=8^{2^j-1}t^{2^j}$.

Weighted-label separation needs a sufficiently large scale $F$, eventually
positive, with $|F|\ge L^t$. The incoming standalone certificate obtains
this from a Pell center $4t+1$ and the divisibility condition
$F=2^{19}t(2u+1)$. That large coefficient dominates its collected quartic.
We move the growth into the Pell center itself.

\subsection{Five new coordinates and a cofinal horizon}

Retain the front coordinates $t,F$ and impose
\begin{equation}\label{eq:newgrowth}
\boxed{\begin{aligned}
 t&=64+4z^2, & S&=t^2, & T&=S^2,\\
 Z&=TF, & X^2-Z^2-2ZF&=1, & F&=T(2u+1).
\end{aligned}}
\end{equation}
The five growth witnesses are $S,T,Z,X,u$; $z$ is the existing horizon
witness. Each displayed residual is quadratic. Eliminating the polynomial
abbreviations in the proof, rather than in the literal quartic, gives
\[
 X^2-\bigl((T+1)^2-1\bigr)F^2=1,\qquad T=t^4.
\]
Thus the center is $a=t^4+1$.

\begin{lemma}[Elementary Pell facts]\label{lem:pell}
For $a\ge2$, define $X_n(a),Y_n(a)$ by
\[
 X_0=1,\quad Y_0=0,\quad
 X_{n+1}=aX_n+(a^2-1)Y_n,\quad Y_{n+1}=X_n+aY_n.
\]
Every integral solution of $X^2-(a^2-1)Y^2=1$ with $Y\ne0$ satisfies
$(|X|,|Y|)=(X_n(a),Y_n(a))$ for a unique $n\ge1$. Moreover
\[
 Y_n(a)\equiv n\pmod{a-1},\qquad
 Y_n(a)\ge(2a-1)^{n-1}.
\]
If $a$ is odd, $\vp(Y_n(a))=\vp(n)$.
\end{lemma}
\begin{proof}
For positive $X,Y$ with $Y>1$, multiply the Pell unit by
$a-\sqrt{a^2-1}$. The new ordinate $aY-X$ is positive because $X<aY$,
and smaller than $Y$ because $X>(a-1)Y$. The new abscissa
$aX-(a^2-1)Y$ is positive by squaring the two positive sides. Descent
ends at $(a,1)$, proving classification and uniqueness. The index
congruence follows by evaluating the integer recurrence at $a=1$.
The second-order recurrence and monotonicity give
$Y_{n+1}=2aY_n-Y_{n-1}\ge(2a-1)Y_n$. For odd $a$, every $X_n$ is odd
and $Y_n$ has the parity of $n$. The duplication identity
$Y_{2n}=2X_nY_n$ then proves the valuation formula.
\end{proof}

\begin{theorem}[Growth without a large fixed multiplier]\label{thm:growth}
Every integer solution of \eqref{eq:newgrowth} has even nonzero $F$ and
\[
 |F|>\bigl(8^{2^{19}-1}t^{2^{19}}\bigr)^t=L^t.
\]
In particular $|F|>t$.
\end{theorem}
\begin{proof}
We have $t\ge64$, $4\mid t$, and $T=t^4>0$.
The odd factor in $F=T(2u+1)$ excludes $F=0$ and makes $F$ even.
Apply Lemma~\ref{lem:pell} at $a=T+1$: $|F|=Y_n(a)$ for $n\ge1$.
Since $T\mid F$ and $Y_n(a)\equiv n\pmod T$, the positive index
satisfies $T\mid n$, hence $n\ge T=t^4$. Therefore
\[
 |F|\ge(2t^4+1)^{t^4-1}>(2t^4)^{t^4-1}.
\]
Put $d=2^{19}$. Since $t^3\ge64^3=d/2$,
$t^4-1>t(d/2-1)$. Also
\[
 (2t^4)^{d/2-1}=2^{d/2-1}t^{2d-4}
   >2^{3d-3}t^d=L.
\]
For the last inequality it suffices that
$t^{d-4}>2^{(5/2)d-2}$; already $t\ge16$ gives the stronger lower
bound $t^{d-4}\ge2^{4d-16}$. Raising to $t$ proves the claim.
The argument makes no assumption about the sign of $F$.
\end{proof}

\subsection{Completeness: the residue condition must survive}

A large Pell solution alone would be insufficient. The subsequent selector
needs $B=\binom{F+t}{t}$ odd, so the completeness construction must select
a large ordinate with a compatible binary residue.

\begin{lemma}[Pell permutation modulo powers of two]\label{lem:permutation}
If $a\equiv1\pmod4$, then $s\mapsto Y_s(a)$ induces a permutation
modulo $2^r$ for every $r\ge1$. Every residue class is attained at
arbitrarily large nonnegative indices; odd residues have odd indices.
\end{lemma}
\begin{proof}
Modulo two, $Y_s\equiv s$. All $X_s$ are $1\pmod4$.
For $h=2^r$, the valuation formula gives $\vp(Y_h)=r$.
The Pell norm and $8\mid a^2-1$ imply
$X_h\equiv1\pmod{2^{2r+2}}$, because $X_h+1$ has valuation one.
The addition formula then gives
\[
 Y_{s+2^r}-Y_s=(X_{2^r}-1)Y_s+X_sY_{2^r}
       \equiv2^r\pmod{2^{r+1}}.
\]
The two lifts of each index class give the two distinct residue lifts.
Induction proves bijectivity. The same identity gives periodicity at each
modulus, and adding periods gives arbitrarily large indices.
\end{proof}

Choose $z$ divisible by $2^{12}$, sufficiently large to accommodate a
given finite refutation and its padded arithmetic trace. Such horizons
are cofinal, and
\[
 t\equiv64\pmod{2^{26}},\qquad \vp(t)=6,
 \qquad \vp(T)=24.
\]
In particular the bit of $t$ at position 24 is zero. Set
\[
 a=T+1,\quad n_0=T,\quad y_0=Y_T(a),\quad a_0=X_T(a),
 \quad b_0=y_0/T.
\]
The index congruence makes $b_0$ integral, and the valuation formula makes
it odd. Since $T$ is even, the duplication formula gives $a_0\equiv1\pmod4$.
Choose $\ell>26$ with $t<2^\ell$. Both $T/2^{24}$ and $b_0$ are
invertible modulo $2^{\ell-24}$. Lemma~\ref{lem:permutation} supplies an
arbitrarily large positive odd $s$ satisfying
\[
 (T/2^{24})b_0Y_s(a_0)\equiv1\pmod{2^{\ell-24}}.
\]
Define
\[
 F=Y_{Ts}(a)=y_0Y_s(a_0),\quad X=X_{Ts}(a),\quad Z=TF,
 \quad u=(F/T-1)/2.
\]
Then $u\in\Z$, every growth row holds, and
\begin{equation}\label{eq:targetresidue}
 F\equiv2^{24}\pmod{2^\ell}.
\end{equation}
The binary supports of $F$ and $t$ are disjoint: below $\ell$, the only
nonzero bit of $F$ is 24, and $t$ has no bits at or above $\ell$.
Consequently addition has no carries. The identity
\[
 \vp\binom{F+t}{t}=s_2(F)+s_2(t)-s_2(F+t)
\]
shows that $B$ is odd.

The remaining canonical construction now applies. Choose an arbitrarily
large odd multiplier containing every odd prime at most $t$ and set
$y=F(2v+1)$. Then every prime at most $t$ divides $y$, so the beta moduli
$1+iy$ for $1\le i\le t+1$ are pairwise coprime. This requires prime
support, not the stronger assertion $t!\mid y$. Finite CRT supplies the
beta code $x$. With $N=B^x$ odd, choose the independent selector
representatives first, and choose the even anchor $C$ last, larger than
every other actual exponent. Its retained congruence gives
$B\mid C+F+1$. Set the actual radix $D=2(C+F+1)$ and $q=D^x$;
then $N\mid q$ and $D>2C$. The inherited digit estimates and main Pell
frame complete the remaining witnesses.

\subsection{The horizon change and the signed bootstrap}

Changing $t=16+4z^2$ to $t=64+4z^2$ changes the allowed horizons. It is
a semantic cofinal replacement, not a coordinatewise map between every
old and new zero. The emitted source audit verifies that the horizon
helper occurs only in its defining row. The retained arithmetic uses
$t\ge16$, $4\mid t$, and the ability to choose a sufficiently large
horizon. All are preserved, and the residue construction above establishes
the needed new cofinal completeness.

Soundness retains the incoming delayed sign proof. The new growth block
first supplies only $|F|>t$ and even nonzero $F$. The even anchor and
even $F$ make $D=2(C+F+1)$ nonzero. The signed exact-power request
$q=D^x$, followed by the beta endpoints, forces $x\ge1$.
The request at exponent $y+1$ then forces $y>0$; the endpoints further
give $x\ge2+(t+1)y$. Now all bases $q+1$ have magnitude at least three.
Their retained requested exponent $F+t$ must be nonnegative, and
$|F|>t$ excludes $F<0$. Thus $F>0$, and the full separation bound is
available. The anchor becomes positive as well. The old radix and binomial
decoding therefore proceed in the original noncircular order.

\subsection{Literal graft and reversible row reduction}

The new raw equation is obtained from the pinned principal210 polynomial
by changing exactly four old rows and appending two rows. The source
coordinates are
\[
 t=x_0,\quad F=x_2,\quad z=x_{178},\quad X=x_{207},
 \quad Z=x_{208},\quad u=x_{209},\quad S=x_{210},\quad T=x_{211}.
\]
The six affected rows are
\begin{align*}
 r_{125}&=x_0-64-4x_{178}^2,\\
 r_{147}&=x_{208}-x_2x_{211},\\
 r_{148}&=x_{207}^2-x_{208}^2-2x_2x_{208}-1,\\
 r_{149}&=x_2-x_{211}-2x_{209}x_{211},\\
 r_{150}&=x_{210}-x_0^2,\\
 r_{151}&=x_{211}-x_{210}^2.
\end{align*}
The other 146 rows are retained literally. This gives
$P_{212}^{\rm sparse}=\sum r_j^2$.

The small-coefficient version is obtained by 465 elementary integral row
operations $r_i\leftarrow r_i+kr_j$, $i\ne j$. Reverse the order and
negate $k$ to invert them. Therefore the transformed rows vanish
simultaneously exactly when the original rows do. Their sums of squares
have exactly the same integer zero set. No variable is added or removed
in this step, and no sign or divisibility condition is introduced.

The independent checker collects the complete quartics and obtains
\[
\begin{array}{c|rrr}
&\text{monomials}&\text{height}&\text{largest residual coefficient}\\\hline
P_{212}^{\rm sparse}&5,290&277,420&226\\
P_{212}^{\rm small}&21,556&6,432&60.
\end{array}
\]
For example the collected coefficient of $x_8x_{10}x_{23}x_{31}$ in
$P_{212}^{\rm small}$ is 6,432. It both witnesses the height and proves
that the degree is exactly four. The denser support is an explicit cost
of the smaller coefficients.

\section{One certificate across two arithmetic stages}
\label{sec:reciprocal-guard}

The direct source has two distinct uses for nonnegativity. A Pell cap
bounds all classified indices before their signs are recovered. A common
digit budget selects the canonical representatives of ten binomial
remainders. Both were represented by the universal ternary polynomial
\[
 \mathcal G(a,b,c)=a^2+b^2+c(c+1).
\]
For integer inputs this is nonnegative, and it represents every
nonnegative integer. Indeed, the three-square theorem applied to $4n+1$
gives one odd and two even roots; division by four gives the displayed
representation.

The new construction shares the certificate across these two stages.
It uses the stronger budget to recover both original inequalities before
either arithmetic stage is interpreted.

\subsection{A general lemma}

Let $U,V\ge0$, $E\ge1$, and $R\ge0$ be integers, and put
\[
 H=UV-UE-VR-1.
\]
If $H\ge0$, then $U,V>0$: setting either to zero makes $H<0$.
The identities
\[
 U(V-E)=VR+1+H>0,\qquad
 V(U-R)=UE+1+H>0
\]
give $V>E$ and $U>R$. Thus
\[
 V-E\ge0,\qquad U-R-1\ge0.                 \tag{RG1}
\]
Conversely, with positive $U,V$,
\[
 \frac{H}{UV}=1-\frac EV-\frac RU-\frac1{UV}. \tag{RG2}
\]
Consequently one equation $H=\mathcal G(a,b,c)$ can replace two separate
ternary certificates for (RG1), if the completeness construction satisfies
the joint bound in (RG2). The $-1$ excludes the simultaneous zero-scale
case, so it must be retained.

\subsection{The literal application}

Here the source subscripts are those of the frozen $219$-coordinate
arithmetic source retained in the archive. Define
\[
 C=2x_{16},\quad Q=x_{126},\quad Y=x_{184},\quad
 U=Q^2,\quad V=Y^2,
\]
and let
\begin{align*}
 (y_i)_{i=1}^{13}={}&(x_{142},x_{151},x_{154}-1,x_{157},x_{160},x_{163},
 x_{166},x_{169},x_{172},x_{175},x_{178},x_{181},x_{187}),\\
 (r_j)_{j=1}^{10}={}&(x_{104},x_{108},x_{111},x_{114},x_{117},
 x_{121},x_{124},x_{128},x_{132},x_{135}),\\
 E={}&C^2+1+\sum_{i=1}^{13}y_i^2,
 \qquad R=4\sum_{j=1}^{10}r_j^2.
\end{align*}
The translated ordinate $x_{154}-1$ is essential. Exact coefficient
comparison verifies that source rows $150$ and $151$ are
\[
 V-E-\mathcal G(x_{207},x_{208},x_{209}),\qquad
 U-R-1-\mathcal G(x_{210},x_{211},x_{212}).
\]
All six guard coordinates occur only in these two rows. Delete
$x_{210},x_{211},x_{212}$, reuse the other triple, and replace both rows by
\[
 \boxed{UV-UE-VR-1-\mathcal G(x_{207},x_{208},x_{209})=0.} \tag{RG3}
\]
Every other residual is retained. The finite source-coordinate
reconstructions supplied by the parent are substituted into these
expressions; their gates are abbreviations, not additional variables.

\subsection{Soundness}

The nonnegativity premises of the general lemma follow immediately from
the displayed sums of squares, without Pell classification or digit
decoding. A zero of (RG3) therefore supplies both old inequalities.
Represent their nonnegative slacks by two fresh ternary triples and
restore the old six guard coordinates. All other rows remain valid,
because those coordinates are private. The certified acyclic unit-pivot
reconstruction restores the parent's eliminated coordinates. Its Pell
cap and subsequent sign bootstrap can then be invoked in their original
order. This proves that every new zero yields a finite syntactic
refutation under the identified inherited calibration.

\subsection{Completeness with a quantitative margin}

Starting from a finite refutation, the inherited canonical construction
chooses a positive even anchor $C$ strictly larger than every other
principal index $k_i$, and then a positive main Pell center $A$ with
\[
 W\ge258,\qquad A>W^W,\qquad C^2<W.
\]
It supplies $Y=Y_C(A)$, $y_i=Y_{k_i}(A)$, $Q=q^x\ge2$, and the ten
normalized remainders $0\le r_j<Q/32$. Write $Z=Y_{C-1}(A)\ge1$.
The recurrence and monotonicity give
\[
 Y_C(A)\ge(2A-1)Z,\qquad y_i\le Z.
\]
Therefore
\[
 \frac EV\le\frac{W+13Z^2}{(2A-1)^2Z^2}
 \le\frac{W+13}{(2A-1)^2}<\frac12,
 \qquad \frac RU<\frac{40}{32^2}=\frac5{128}.
\]
The first strict inequality follows already from $A>W\ge258$.
Also $C\ge2$, so $Y_C(A)\ge Y_2(A)=2A\ge4$. Hence $U\ge4$,
$V\ge16$, and
\[
 \frac EV+\frac RU+\frac1{UV}
 <\frac12+\frac5{128}+\frac1{64}
 =\frac{71}{128}<1.
\]
Thus $H>0$. Its ternary representation supplies the new triple. The
CRT representatives and powers are chosen first, the main Pell frame
after them, and this final representation last. There is no circular
witness choice.

This is an equivalence of existence of zeros, established by
canonical-witness completeness. It is not a pointwise bijection between
every old zero and every new zero.

\subsection{Composing the linear orientation numerator}

The independent orientation improvement replaces
\[
 \Omega_0=\sum_i\bigl((X_i-1)^2+(Y_i-k_i)^2\bigr)
 \quad\hbox{by}\quad
 \Omega_1=\sum_i\bigl(1-X_i+(Y_i-k_i)^2\bigr).
\]
Under the bounded raw Pell classification,
$(X_i,Y_i)=\epsilon_i(X_{\nu_i}(A),Y_{\nu_i}(A))$, reduction modulo
$A-1$ gives the nonnegative error
\[
 \sum_i\bigl(1-\epsilon_i+(\epsilon_i\nu_i-k_i)^2\bigr).
\]
It vanishes precisely when every $\epsilon_i=1$ and $\nu_i=k_i$.
The same signed rational-integrality argument authenticates this error
simultaneously with all power requests. The complete bound and
restoration proof are retained in the accompanying orientation proof.
The two changes concern disjoint private rows and therefore compose.

The recommended resulting polynomial has $82$ unrestricted integer
variables, $23$ residual squares, exact degree $16\,964$, and $1\,723$
live binary arithmetic operations. Its lower-dimensional companion has
$81$ variables, $22$ residual squares, degree $2\,097\,732$, and
$1\,711$ operations. A lower-degree companion has $106$ variables,
degree $32$, $47$ residual squares, and $1\,830$ operations. All three
retain maximum absolute literal $1\,099\,511\,627\,780$; this is not
the expanded coefficient height. Conservative height bounds remain
$2^{59\,805}$, $2^{7\,342\,557}$, and $2^{308}$, respectively.

Exact sparse identities bind the source predicates, structural
substitution checks bind all retained expressions and unit pivots, and
nonzero top homogeneous coefficients establish the stated degrees.
Nonzero derivative certificates establish genuine dependence on every
input. These finite certificates complement the unbounded proof; the
consistency equivalence still uses the explicitly identified inherited
C222/ZF/ZFC calibration.


\subsection{Full bounded-error argument for the linear orientation}

Here is the unbounded estimate underlying the orientation substitution.
It applies at the raw-classification stage, before requested indices have
been authenticated. Put $D=A-1$, assume $|A|>W^W$, $W\ge258$, and write
$b=\lfloor\sqrt{W-1}\rfloor$. The retained anchor and comparison rows
classify at most 15 principal pairs as
\[
 (X_i,Y_i)=\epsilon_i(X_{\nu_i}(A),Y_{\nu_i}(A)),
 \quad\epsilon_i\in\{1,-1\},\quad |\nu_i|,|k_i|\le b.
\]
Signed-index Pell polynomials satisfy $X_n(A)\equiv1$ and
$Y_n(A)\equiv n\pmod D$. Therefore
\[
 \Omega_1\equiv o:=\sum_i
       \bigl(1-\epsilon_i+(\epsilon_i\nu_i-k_i)^2\bigr)\pmod D.
\]
Every summand of $o$ is nonnegative, and it vanishes exactly at
$\epsilon_i=1$, $\nu_i=k_i$. Moreover $0\le o\le90b^2$.
In the standalone orientation certificate this bound already proves that
$D\mid\Omega_1$ is equivalent to the desired canonical orientation.
Notice that $\Omega_1$ itself need not be nonnegative; its bounded residue
$o$ has the required sign.

For joint power authentication, let $n_j$ be the at most five bases,
and let $(m_i,k_i)$ be their at most seventeen requests. Every base,
requested output, and index has absolute value at most $b$. Define
\begin{align*}
 T_j&=2An_j-n_j^2-1,\\
 E_j&=\sum_{i\in j}\bigl((m_i-X_i+AY_i)^2-Y_i^2
                    -2n_jY_i(m_i-X_i)\bigr).
\end{align*}
The literal joint residual is
\begin{equation}\label{eq:jointlinear}
 \Omega_1\prod_jT_j
 +D\sum_j n_jE_j\prod_{\ell\ne j}T_\ell
 -zD\prod_jT_j=0.
\end{equation}
All denominators in its mathematical interpretation are nonzero. A zero
base gives $T_j=-1$; otherwise $|T_j|>W^W$ and $n_j/T_j$ has the sign
of $A$. Also $D$ has the sign of $A$ and $|D|\ge W^W$.

For clarity, the relevant power congruence follows from the identity
\[
 (m-X+AY)^2-Y^2-2nY(m-X)
   =(m-X+(A-n)Y)^2+(2An-n^2-1)Y^2.
\]
Modulo $T=2An-n^2-1$, the expression $X_{\nu}(A)+(n-A)Y_{\nu}(A)$
is $n^{\nu}$; this follows from its recurrence, starting at $1,n$.
For negative indices this is a modular inverse, permitted because
$\gcd(n,T)=1$. Equivalently one may clear all powers before stating the
congruence.

Use the latter form by taking the mathematical clearing integer
\[
 B_* =\prod_{j:n_j\ne0}n_j^b,\qquad
 \delta_i=B_*m_i-\epsilon_i(B_*/n_j^b)n_j^{b+\nu_i}.
\]
These are integers since $b+\nu_i\ge0$. The congruences and
\eqref{eq:jointlinear} imply that
\[
 \frac{B_*^2o}{D}
   +\sum_{j:n_j\ne0}\sum_{i\in j}\frac{n_j\delta_i^2}{T_j}
\]
is an integer. Every nonzero summand has the same sign, that of $A$.
Since $|B_*|\le b^{5b}$ and $|\delta_i|\le2b^{6b}$, its total absolute
numerator estimate is
\[
 B_*^2o+\sum_i|n_j|\delta_i^2
 \le90b^{10b+2}+68b^{12b+1}
 \le158b^{12b+2}<W^{6b+2}\le W^W.
\]
Here $b\ge16$, so $158<b^2$, and $6b+2\le W$. Dividing by the lower
denominator bound shows that the displayed integer has absolute value
strictly below one. It must be zero. Its same-sign terms vanish
separately, giving $o=0$ and every $\delta_i=0$. Thus the indices and
signs are canonical and every nonzero-base request is the exact signed
power. In particular a base of absolute value at least two cannot have
a negative requested exponent with an integral output.

The source's zero-base bootstrap is retained until the front equations
exclude those temporary cases. Conversely, canonical pairs make
$D\mid\Omega_1$, and exact powers make $T_j\mid E_j$ for nonzero bases;
zero-base denominators are units. Their quotient sum supplies the new $z$.
This proves completeness without a pairwise-coprimality assumption on the
different $T_j$. The entire argument uses bounded classified indices
before canonicality, so it does not assume its own conclusion.

\section{Two Pell roots carry two quotients}
\label{sec:two-pell-carriers}

This section gives an independent compact improvement, retaining the exact
admissible942 coefficient tuple of the incoming \emph{Pell Carrier65}
package.  That package removes the first-index quotient by using the first
Pell root.  We now use a second Pell root to remove the native transport
quotient.  The factor four in an odd-ordinate congruence is what makes the
second reduction work without a parity restriction on the auxiliary base.
The resulting ordinary polynomial has thirteen strictly positive variables,
69 live binary arithmetic operations, and exact degree~174.  Its standard
unrestricted-integer presentation has~39 variables, 160 operations, and
degree~348.  A companion trades one extra operation for degrees170 and340.

The parent and its current calibration are retained under the compact
component's \texttt{inherited/} directory.  In particular the complete
admissible942 proof, rather than an arbitrary four-number tuple satisfying
elementary inequalities, gives these equations their consistency meaning.
The general use of a Pell ordinate as a quotient carrier already occurs
in package~170.  The present contribution is the stronger residue, its
use for native sign recovery before rank decoding, and its composition
with the first-root construction.

\subsection{The entire polynomial}

There are exactly thirteen independent positive integer inputs,
in the order used by the machine-readable equation:
\[
 J,F,\alpha,f,i,T,s,w,n,\eta,\zeta,m,\gamma>0.
\]
Retain the four fixed integers $\beta,K_0,M_C,M_F$ from the exact parent.
Their unsigned big-endian byte strings total~94,689 bytes, and their
largest bit length is~189,693.  The compact component includes their complete
finite reconstruction, its~399-character seed, and all four literal
byte strings.  The polynomial below contains no variable exponent,
division, floor, factorial, or unspecified predicate.  Every other letter
is a polynomial abbreviation:
\begin{align*}
 h&=\beta J,& q&=h+1,& C&=q-F-\alpha,\\
 X&=wq,&Y&=sq^3,&E&=XY,\\
 k&=\eta+\zeta,&\ell&=EkY,&c&=kY+\eta,\\
 a&=E+Y,&H&=4a+3,&\Delta&=a^2+H=(a+2)^2-1,\\
 D&=X+ac+\gamma H,&\tau&=En+k-R,&B&=(K_0+w)C-F,
\end{align*}
where
\[
 R=[h(q-F)+\alpha](q^2-1)+(M_C+qM_F)J.
\]
The last definitions are
\begin{align*}
 b&=cT-R,& V&=f^2b-c,&S&=hi\Delta c^2,\\
 L&=V^2,& y&=1+4\bigl[h(m-LB)-B\bigr],\\
 N_0&=\tau^2-\ell(\ell+k),&
 N_1&=D^2-\Delta c^2,\\
 N_a&=S^2(V^2-y^2)+y^2,&
 N_s&=\Delta f^2-S^2.
\end{align*}
The single equation is
\begin{equation}
 \boxed{P=N_0N_1N_aN_s-\Delta=0.}
 \label{eq:second-carrier}
\end{equation}
The companion changes only $L$ to $bV$.  Notice that $y$ is a signed
polynomial value, not an additional positive unknown.

\subsection{A residue with a factor of four}

For an integer $S\ge2$, write
\[
 \chi_S(r)+\psi_S(r)\sqrt{S^2-1}
   =(S+\sqrt{S^2-1})^r.
\]
The odd subsequence $u_j=\psi_S(2j+1)$ satisfies
\[
 u_0=1,\quad u_1=4S^2-1,\quad
 u_{j+2}=(4S^2-2)u_{j+1}-u_j.
\]
Induction modulo $4S^2$ therefore proves
\begin{equation}
 \psi_S(2j+1)\equiv(-1)^j\pmod{4S^2}.
 \label{eq:odd-four-residue}
\end{equation}
This holds for odd as well as even $S$.  The weaker congruence modulo
$S^2$ would be insufficient for the sign selection used below.

Suppose that integers $V,y$ satisfy
\begin{equation}
 S^2V^2-(S^2-1)y^2=1.\label{eq:second-adjacent}
\end{equation}
Neither coordinate is zero.  The positive solutions of
$x^2-(S^2-1)u^2=1$ are exactly
$(\chi_S(r),\psi_S(r))$, $r\ge1$.  One can prove this directly by
descent: for $u>1$, multiplication by $S-\sqrt{S^2-1}$ sends the
ordinate to $Su-x$, which is positive and smaller than $u$, because
$S-1<\sqrt{S^2-1}<x/u<S$.  The new abscissa is positive as well.
The terminal positive pair is $(S,1)$, and reversing the descent gives
the stated powers.

Apply this to $x=S|V|$ and $u=|y|$.  Modulo $S$, the even-index
$\chi$ values are $+1$ or $-1$, whereas the odd-index values are zero.
The index is consequently odd.  Equation\eqref{eq:odd-four-residue}
then gives
\[
 y\equiv+1\text{ or }-1\pmod{4S^2}.
\]
If $y\equiv1\pmod4$, the minus sign is impossible.  Thus
\begin{equation}
 y\equiv1\pmod4\quad\Longrightarrow\quad4S^2\mid y-1.
 \label{eq:carrier-full-divisibility}
\end{equation}
Also, directly from\eqref{eq:second-adjacent},
\begin{equation}
 |y|<2|V|.\label{eq:carrier-root-bound}
\end{equation}
No packed-rank interpretation or positivity of $V$ or $y$ was used.

\subsection{Soundness before rank recovery}

On the entire positive domain, $h\ge31$, $c>1$, $\Delta>1$ and
$S>c$.  At a zero of\eqref{eq:second-carrier}, cancellation of
the positive integer $\Delta$ gives
\[
 N_0N_1N_a\bigl[f^2-\Delta(hi c^2)^2\bigr]=1.
\]
All four factors are integer units.  Each negative sign is excluded
independently.  The consecutive-product negative-Pell obstruction excludes
$N_0=-1$, since
$\ell(\ell+k)=EY(EY+1)k^2$ and $EY\ge2$.  For completeness, if
$x^2-t(t+1)u^2=-1$ with $t\ge2$ and least positive $u$, then
$tu<x<(t+\tfrac12)u$ after taking absolute values; multiplication by
$(2t+1)-2\sqrt{t(t+1)}$ gives a positive smaller ordinate
$(2t+1)u-2x$, a contradiction.  This does not assume $\tau>0$.

Because $\Delta$ is zero or three modulo four, neither $N_1$ nor the
normalized strong norm can be $-1$.  Finally $N_a$ is congruent modulo
four to either $y^2$ or $V^2$, and cannot be $-1$.  Hence
\begin{equation}
 N_0=N_1=N_a=1,\qquad N_s=\Delta.
 \label{eq:four-unit-signs}
\end{equation}
In particular $N_a=1$ is\eqref{eq:second-adjacent}.  Our definition
of $y$ gives $y\equiv1\pmod4$, so
\eqref{eq:carrier-full-divisibility}, together with $h\mid S$, yields
\[
 h^2\mid h(m-LB)-B,\qquad\text{and therefore }h\mid B.
\]
No primality or parity assumption on $h$ enters this deduction.

For the main profile, $L=V^2\ge|V|\ge1$.  For the companion
$L=bV$, the same inequality holds at a zero.  Indeed $b=0$ would give
$V=-c$ and
\[
 y^2-c^2=\frac{c^2-1}{S^2-1}\in(0,1),
\]
impossible for integers.  Equation\eqref{eq:four-unit-signs} gives
$f^2=1+\Delta h^2i^2c^4>c$, so $V=f^2b-c$ has the sign of nonzero
$b$.  Therefore $bV=|b||V|\ge|V|$.

Set $W=1+hL>0$.  If $B\le-1$, positivity of $m$ forces
\[
 y=1+4(hm-WB)\ge1+4(h+W)>4hL\ge4|V|,
\]
contradicting\eqref{eq:carrier-root-bound}.  Thus $B\ge0$.
The omitted native quotient is consequently recovered as the positive
integer
\[
 z=B/h+1.
\]
Its parent factor is exactly
\[
 N_t=(K_0+w)C+(q-F)-zh=B+q-zh=1.
\]
Keep the other parent inputs, replace its strong multiplier by $hi$,
and replace its positive auxiliary ordinate by $|y|$.  Every parent
factor now has the required value, and the Carrier65 polynomial vanishes.
This is the complete soundness map.  It uses no conclusion about $C$,
$R$, or the first Pell index before reaching the already proved parent.

\subsection{Completeness with a fresh strong seed}

A finite syntactic ZFC refutation supplies the parent's canonical outer
data.  Among their properties are
\[
 N_0=N_1=1,\quad R\equiv3\pmod4,\quad
 c=\psi_{a+2}(R),\quad X=2^R,
\]
as well as $C>0$, $F<q$, $h\mid B$, and $R>3q$.  The inherited
first-root coordinate $n$ is positive.  These data are independent of
the private strong seed and the auxiliary coordinates.

There is actually a strict native surplus: $2^R>q^2$ gives $w=X/q>q$,
and hence
\begin{equation}
 B=(K_0+w)C-F>w-q>0.\label{eq:strict-native-surplus}
\end{equation}
Choose a positive solution
\[
 f^2-\Delta I^2=1,\qquad hc^2\mid I.
\]
Such solutions exist cofinally.  In the finite ring
$(\mathbb Z/(hc^2)\mathbb Z)[t]/(t^2-\Delta)$, the element $a+2+t$
is a unit of norm one and therefore has finite order.  Any positive
multiple of that order gives $I\equiv0\pmod{hc^2}$; its positive
powers give arbitrarily large choices.  Set $i=I/(hc^2)>0$ and
$S=\Delta I$.  Then $N_s=\Delta$ for our definitions.

Put
\[
 V_*=\frac{\chi_S(R)}{S},\qquad y_* =\psi_S(R).
\]
The quotient is a positive integer because $R$ is odd.  Write $R=2j+1$
and $\chi_S(R)/S=Q_j(S^2)$.  The elementary identities retained in
the parent proof are
\[
 Q_j(0)=(-1)^jR,\qquad
 Q_j(1-A^2)=(-1)^j\psi_A(R),\quad A=a+2.
\]
Since $j$ is odd, $S^2=1-A^2\pmod{f^2}$ and $c\mid S$, they give
$V_*\equiv-c\pmod{f^2}$ and $V_*\equiv-R\pmod c$.
Moreover $\gcd(c,f)=1$ and $f^2\equiv1\pmod c$.  Thus
\[
 T=\frac{V_*+c+Rf^2}{cf^2}>0
\]
is an integer and our polynomial value $V$ is $V_*$.  These quotients
construct witnesses; none is an operation in the delivered polynomial.

Equation\eqref{eq:odd-four-residue} gives
$y_*\equiv-1\pmod{4S^2}$.  Use the signed value $y=-y_*$, so
$y\equiv1\pmod{4S^2}$, and define
\begin{equation}
 m=\frac{y-1}{4h}+\frac{(1+hL)B}{h}.
 \label{eq:new-positive-carrier-witness}
\end{equation}
It is an integer.  It is positive because $B$ is a positive multiple of
$h$ by\eqref{eq:strict-native-surplus}, $L\ge V\ge1$, and
$|y|<2V$:
\[
 m> -\frac{2V+1}{4h}+1+hL>0.
\]
Equation\eqref{eq:new-positive-carrier-witness} makes the literal
carrier equal the chosen signed $y$, and all four factors in
\eqref{eq:four-unit-signs} have their required values.  Thus $P=0$.
This is an existence construction that refreshes the private Pell seed;
it is not a bijection on every old and new zero.

\subsection{Exact size and unrestricted integer variables}

Relative to the~65-operation parent, remove the three final native-factor
rows, add the one subtraction defining $B$, remove one product
multiplication, scale the strong multiplier by $h$ using one multiplication,
and compute $y$ with six operations.  The count is therefore
$65-3+1-1+1+6=69$.  The companion requires one extra multiplication for
$bV$.  Squares and scalar multiplications are charged as binary operations.

The respective factor degrees are
\[
\begin{array}{c|rrrr|r}
 &N_0&N_1&N_a&N_s&P\\\hline
L=V^2&22&18&86&48&174\\
L=bV&22&18&82&48&170
\end{array}
\]
The main-norm degree uses the exact cancellation, with $Z=X+\gamma H$,
\[
 (ac+Z)^2-(a^2+H)c^2=2acZ+Z^2-Hc^2.
\]
Let $q_0=\beta J$ and $k_0=\eta+\zeta$.  The complete leading
homogeneous form of the main equation is
\[
 -128q_0^{99}\gamma i^4k_0^{15}w^{14}s^{27}f^8T^4
          (q_0-F-\alpha)^2.
\]
The companion replaces $f^8$ by $f^4$.  The coefficient selecting
$\eta^{15}$ and $F^2$ is $-128\beta^{99}$, which is nonzero on the
actual fixed tuple.  These are exact degrees, not only graph estimates.

To obtain unrestricted integer equations, replace each positive input by
\[
 1+u^2+v^2+w(w+1),\qquad u,v,w\in\mathbb Z.
\]
The classical three-square theorem applied to $4N+1$ proves that every
$N\ge0$ is $u^2+v^2+w(w+1)$: exactly one of the three squares is odd,
and subtracting its1 and dividing by4 gives the representation.
Conversely this expression is nonnegative on all integer triples.
The substitution uses seven operations per coordinate.  It gives
$39$ integer variables and respectively $160$ or $161$ operations.
Degrees double to348 and340.  Nonzero highest homogeneous terms survive:
setting each second and third new coordinate to zero doubles the
monomial exponents injectively.

The independent compact checker verifies~25 or~26 exact formula cuts,
the full~96-term leading forms, all coordinate and coefficient bindings,
the actual fixed-tuple degree witnesses, four coefficient reconstructions,
and~1,278 finite Pell-carrier fixtures.  Targeted changed gates are rejected.
These finite checks complement the preceding unbounded proof.  They do
not replace the inherited native/K215/refutation calibration or provide a
fresh end-to-end proof-assistant derivation.  Composing that calibration
with the two witness maps and the exact domain translation proves in ZFC
that the integer equation has no zero if and only if
$\operatorname{Con}(\mathrm{ZFC})$.

The result is a tradeoff against the matched parent and inspected compact
family.  Package~170 has a~13-positive-variable, 70-operation degree~206
point and a~71-operation degree~162 point, with a different~295,388-byte
tuple.  The present~69-operation degree~174 point does not assert a
minimum over all Diophantine equations, all integer domains, or all
coefficient architectures.


\section{Reproduction, evidence, and integration}

\subsection{Files that define the equations}

The archive is organized around three independently usable components:
\begin{itemize}
\item \pathtext{quartic212/} contains both literal 212-variable systems,
their collected quartics, the frozen principal210 input, the source-change
certificate, the reversible row ledger, builders, evaluator, and an
independent checker.
\item \pathtext{projected/} contains the 106-, 82-, and 81-variable direct
graphs, their exact input and reconstruction maps, frozen parent graphs,
the immediate parent archive, and separate construction and verification
programs.
\item \pathtext{compact13/} contains the two 13-positive-variable graphs,
their 39-integer lifts, the fixed coefficient tuple and its reconstruction,
the exact incoming Carrier65, and ordinary proofs and checkers.
\end{itemize}
The top-level \pathtext{verify.py} runs the supplied finite checks. It does
not search for roots. The checksum manifest authenticates the delivered
bytes. The report source and all included equation definitions are editable.

\subsection{What was actually checked}

For the quartic, a checker independent of the producer verifies the pinned
source hash, all 146 unchanged residuals, the six displayed affected rows,
all 212 live coordinates, and the horizon helper's sole occurrence. It
replays all 465 row operations forward and backward and independently
collects both polynomials. Its seven targeted corruption checks reject
changed growth coefficients and a changed row operation. Pell recurrence
fixtures are auxiliary tests only.

For the projected equations, source predicate identities identify the two
old guards, structural substitutions identify every retained residual,
and acyclic unit-pivot certificates reconstruct the eliminated coordinates.
The linear-orientation replacement is checked as an exact residual change,
with a separate ordinary proof of existential equivalence. Nonzero
homogeneous terms certify the exact degrees; derivative certificates
establish dependence on each listed input. Rebuilding the component
reproduces its new JSON artifacts byte for byte.

For the compact equations, an independent checker verifies the formal
cuts in the evaluator, the full 96-term leading homogeneous forms, the
actual fixed coefficient bytes, and source bindings. The report gives
explicit nonzero leading coefficients after specialization. The
39-integer lifts are checked substitutions of the positive parametrization.
All substantive verifiers retain their checks under Python's \pathtext{-O}
mode; none relies on disabled assertions for correctness.

No finite check asserts that one of these consistency equations has a
root. An evaluation of a test tuple, or a modular recurrence fixture, is
an arithmetic test and does not constitute a ZFC refutation. The
unbounded theorems remain the mathematical arguments given here and in
the named inherited sources.

\subsection{Suggested repository placement}

This package can be incorporated under the repository's
\pathtext{Computability/DiophantineConsistency/} subtree, in its
\pathtext{con-zfc-diophantine-forms/} collection.
The directory name should be assigned by the repository's intake process,
since numbering may have advanced after the pinned snapshot. The supplied
\pathtext{INTEGRATION.md} records the recommended claims and their
qualification. The work creates a reviewable submission and does not
modify the remote repository.

\section{Further research questions}

\begin{enumerate}
\item \textbf{Combine the new growth center with the projected family.}
The quartic construction and the 82-variable construction currently descend
from different exact immediate sources. A fully authenticated composition
could replace the projected family's large growth literals, then rerun its
unit-pivot search. The required work is a source-port and frame audit,
followed by degree certificates. Counts cannot be obtained safely by merely
subtracting the separate savings.

\item \textbf{Reduce the coefficient cost of the quartic without excessive
support growth.} The 465-step row reduction has a clear certificate but
optimizes a quadratic coefficient-vector norm rather than final support
or expanded coefficient height. Search over small integral row operations,
with a lexicographic or Pareto objective, may improve the 212-variable
tradeoff. A further reduction must publish its reversible ledger and
recollected polynomial.

\item \textbf{Generalize the shared-bound certificate.}
The expression $UV-UE-VR-1$ combines two relative deficits because the
canonical completeness construction has substantial spare margin.
For several positive scales $U_i$ and nonnegative deficits $E_i$, clearing
denominators in $1-\sum_i E_i/U_i$ suggests further pooling. The main
challenge is to exclude all zero-scale cases and recover each original
inequality before any dependent decoding. The resulting degree and gate
cost may outweigh the saved guard variables.

\item \textbf{Lower degree after the quotient deletion.}
The 13-positive-variable carrier removes one quotient using a second Pell
ordinate, but its positive shift raises the degree. The $bV$ companion
already reduces degree by four. A different positive multiplier bounded
below by $|V|$, with a lower total degree and no new quantified coordinate,
could improve this point. Soundness must handle the sign of $b$ and the
case $b=0$ before using the native decoder.

\item \textbf{Formalize the new arithmetic lemmas and seal the full bridge.}
The general fused-bound lemma, the odd-Pell congruence modulo $4S^2$, and
the polynomial-center growth lemma are compact formalization targets.
Afterward, bind the emitted coefficient files to a closed theorem about
the fixed object-language ZFC proof predicate. The existing ordinary
calibration and finite checker receipts are useful inputs to that work;
they do not themselves constitute the completed new formal endpoint.
\end{enumerate}

\begin{thebibliography}{9}
\bibitem{repo}
V. Reshetnikov and contributors, \emph{ProveIt}, pinned revision
\pathtext{b42c796d1f5de9912a576bc2fd77b778ae7402bf}.
\href{https://github.com/VladimirReshetnikov/ProveIt/tree/b42c796d1f5de9912a576bc2fd77b778ae7402bf/Computability/DiophantineConsistency}{Repository arithmetic directory}.

\bibitem{growth}
\emph{Pell growth from divisibility: an explicit quartic for Con(ZFC)},
incoming \pathtext{ConZFC_Pell_Growth_210_2026-10-25.zip}, same snapshot.
The exact principal210 JSON has SHA-256
\pathtext{20a2108353a3756ed7d1ddfdc71695131dcf0b6d94039988d9565ac48e6e1f0a}.
Selected ordinary proofs and attribution are reproduced in
\pathtext{quartic212/upstream/}.

\bibitem{direct}
\emph{Direct integer elimination}, incoming
\pathtext{ConZFC_Direct_Integer_Elimination_2026-10-25.zip}, same snapshot.
The exact immediate archive is retained in
\pathtext{projected/sources/parent_archive.zip}.

\bibitem{carrier}
\emph{A first Pell carrier: 65 operations}, incoming
\pathtext{ConZFC_Pell_Carrier65_2026-10-25.zip}, same snapshot.
Its exact admissible942 coefficient tuple and carrier65 input are retained
with the compact component.

\bibitem{pool}
ProveIt packages 135 (cofinal anchors), 151 (factorial-free binomial CRT),
157 (linear CRT), 173 (CRT index anchor), and 176 (oriented pooling).
These supply the inherited arithmetic interfaces cited in the respective
component proof maps. Exact source hashes and retained paths appear in
the provenance files.

\bibitem{calibration}
ProveIt, \pathtext{KEndpoint.lean} and
\pathtext{ConstructibleChoice.lean}, identified calibration endpoints.
See the retained ordinary dependency audit for their pinned URLs, proof
terms, immediate dependencies, and source-review limitations.

\bibitem{legendre}
P. Pollack and P. Schorn, \emph{Dirichlet's proof of the three-square
theorem: an algorithmic perspective}, Mathematics of Computation 88
(2019), no.\ 316, 1007--1019,
\href{https://doi.org/10.1090/mcom/3349}{doi:10.1090/mcom/3349}.
The only representation fact used here is
the standard consequence $n=a^2+b^2+c(c+1)$ for $n\ge0$; its reduction to
three squares is given in the text.
\end{thebibliography}

\clearpage\appendix
\section{The complete small-coefficient quartic}\label{app:literal}

All 212 coordinates $x_0,\ldots,x_{211}$ range independently over $\Z$.
The equation is $\sum_{j=0}^{151}R_j^2=0$. The following literal rows are
generated from \pathtext{quartic212/reduced/equation.json}; a separate
parser verifies every printed coefficient and variable index against that
file. Their semantic names before row reduction are in the variable table
supplied with the equation. Row reduction changes the residual basis but
does not change the coordinate domain.

{\small
\setlength{\jot}{2pt}
\allowdisplaybreaks[2]
% BEGIN RESIDUAL 0
\begin{align*}
R_{0} ={}& x_{0} - x_{0} x_{4} - x_{4} + x_{7}
\end{align*}
% END RESIDUAL 0
% BEGIN RESIDUAL 1
\begin{align*}
R_{1} ={}& -x_{2} - x_{4} x_{5} + x_{6} x_{7}
\end{align*}
% END RESIDUAL 1
% BEGIN RESIDUAL 2
\begin{align*}
R_{2} ={}& -1 + x_{1} - x_{6} x_{7}
\end{align*}
% END RESIDUAL 2
% BEGIN RESIDUAL 3
\begin{align*}
R_{3} ={}& 1 + x_{2} - x_{3} x_{18} + 2x_{16}
\end{align*}
% END RESIDUAL 3
% BEGIN RESIDUAL 4
\begin{align*}
R_{4} ={}& -1 - x_{1} + x_{6} x_{7} - 3x_{8} - x_{8}^{2} \\
&{}  - x_{8} x_{9} - x_{8} x_{10} - 4x_{9} + x_{9} x_{10} + 2x_{10}^{2} \\
&{}  + x_{19} - x_{23} - x_{71}
\end{align*}
% END RESIDUAL 4
% BEGIN RESIDUAL 5
\begin{align*}
R_{5} ={}& 2x_{8} - 2x_{8}^{2} - 8x_{8} x_{20} - 8x_{20} - 8x_{20}^{2} \\
&{}  + x_{21} + x_{33}
\end{align*}
% END RESIDUAL 5
% BEGIN RESIDUAL 6
\begin{align*}
R_{6} ={}& 1 + 2x_{21} - 4x_{21}^{2} + x_{22}
\end{align*}
% END RESIDUAL 6
% BEGIN RESIDUAL 7
\begin{align*}
R_{7} ={}& -1 - x_{1} + x_{6} x_{7} - x_{9}^{2} - 2x_{9} x_{10} \\
&{}  - x_{10}^{2} + x_{23}
\end{align*}
% END RESIDUAL 7
% BEGIN RESIDUAL 8
\begin{align*}
R_{8} ={}& 1 - x_{8}^{2} - 2x_{8} x_{10} - x_{10}^{2} + x_{24}
\end{align*}
% END RESIDUAL 8
% BEGIN RESIDUAL 9
\begin{align*}
R_{9} ={}& 2 + x_{1} - x_{6} x_{7} - 13x_{8} + 8x_{8}^{2} \\
&{}  + 4x_{8} x_{9} - 14x_{8} x_{10} - 16x_{8} x_{19} - 4x_{8} x_{20} - 12x_{8} x_{23} \\
&{}  - 24x_{8} x_{24} + 4x_{8} x_{52} + 4x_{9} - 3x_{9}^{2} + 2x_{9} x_{10} \\
&{}  + 8x_{9} x_{23} - 32x_{9} x_{36} - 38x_{10} - 5x_{10}^{2} + 2x_{10} x_{11} \\
&{}  - 52x_{10} x_{19} + 8x_{10} x_{24} + 4x_{10} x_{27} - 24x_{10} x_{36} + 8x_{10} x_{52} \\
&{}  + x_{11}^{2} - 17x_{19} + 16x_{19} x_{20} - 8x_{19} x_{23} - 24x_{19} x_{24} \\
&{}  - 8x_{19} x_{27} - 6x_{20} - 4x_{20}^{2} - x_{23} + 4x_{23}^{2} \\
&{}  + 16x_{23} x_{24} - 8x_{23} x_{36} - 29x_{24} + 4x_{24}^{2} + x_{25} \\
&{}  - 12x_{25}^{2} + 3x_{26} + 10x_{27} + x_{28} + x_{31} \\
&{}  + x_{33} + 4x_{36} - 4x_{36}^{2} + x_{37} - x_{39} \\
&{}  - x_{42} - 8x_{43} + 4x_{52}^{2} - x_{53} - x_{59} \\
&{}  + 12x_{71}
\end{align*}
% END RESIDUAL 9
% BEGIN RESIDUAL 10
\begin{align*}
R_{10} ={}& -1 - x_{1} + x_{6} x_{7} + 10x_{8} - 2x_{8}^{2} \\
&{}  - 6x_{8} x_{9} - 2x_{8} x_{10} - 8x_{8} x_{23} + x_{9}^{2} + 3x_{10}^{2} \\
&{}  - x_{23} - 4x_{25}^{2} + x_{26} - 2x_{43} - 2x_{71}
\end{align*}
% END RESIDUAL 10
% BEGIN RESIDUAL 11
\begin{align*}
R_{11} ={}& 2x_{8}^{2} + 4x_{8} x_{9} + 3x_{9}^{2} + 2x_{9} x_{10} - 6x_{10}^{2} \\
&{}  + 2x_{10} x_{11} - 16x_{10} x_{23} - 16x_{10} x_{24} + x_{11}^{2} - 3x_{23} \\
&{}  - 4x_{23}^{2} - 8x_{23} x_{24} - 2x_{24} - 4x_{24}^{2} + x_{27} \\
&{}  - x_{39} - 4x_{71}
\end{align*}
% END RESIDUAL 11
% BEGIN RESIDUAL 12
\begin{align*}
R_{12} ={}& -2 - 2x_{1} + x_{2} + x_{4} x_{5} + x_{6} x_{7} \\
&{}  + 3x_{8} + 3x_{8}^{2} + 6x_{8} x_{10} + 4x_{9} + 23x_{10} \\
&{}  + 7x_{10}^{2} - 2x_{10} x_{11} - 16x_{10} x_{19} + 8x_{10} x_{24} + 4x_{10} x_{27} \\
&{}  - x_{11}^{2} + 9x_{19} - 16x_{19} x_{24} - 8x_{19} x_{27} - x_{20} \\
&{}  + x_{23} + 22x_{24} + 12x_{27} + x_{28} + x_{39} \\
&{}  - x_{42} + 2x_{71} - x_{131}^{2} - x_{131} x_{136} - x_{131} x_{161} \\
&{}  - x_{133}^{2} + x_{134}^{2} - x_{135} + x_{137} - x_{160} \\
&{}  + x_{162}
\end{align*}
% END RESIDUAL 12
% BEGIN RESIDUAL 13
\begin{align*}
R_{13} ={}& 4 + x_{0} + 3x_{1} - 2x_{2} - 2x_{2} x_{208} \\
&{}  - x_{3}^{2} + x_{3} x_{18} + 2x_{3} x_{131} - x_{4} x_{5} - 2x_{6} x_{7} \\
&{}  + 10x_{8} + 23x_{8}^{2} + x_{8} x_{9} - 3x_{8} x_{10} + 24x_{8} x_{19} \\
&{}  + 8x_{8} x_{23} - 15x_{9} - 25x_{9}^{2} + 9x_{9} x_{10} - 8x_{9} x_{23} \\
&{}  - 4x_{9} x_{24} + 4x_{9} x_{34} - 51x_{10} - 5x_{10}^{2} + 2x_{10} x_{11} \\
&{}  + 4x_{10} x_{19} - 8x_{10} x_{23} - 16x_{10} x_{27} - 8x_{10} x_{57} - 2x_{11} \\
&{}  - 2x_{16} + 3x_{19} - 8x_{19} x_{23} - 8x_{19} x_{24} - 8x_{19} x_{27} \\
&{}  + 2x_{20} - x_{23} - 8x_{23}^{2} + 16x_{23} x_{24} + 9x_{24} \\
&{}  - 8x_{24}^{2} - 16x_{24} x_{27} - x_{25} - 4x_{25}^{2} + x_{26} \\
&{}  + 10x_{27} - 4x_{27}^{2} + x_{28} + x_{29} + x_{31} \\
&{}  + 4x_{34}^{2} - x_{35} - x_{36} - x_{39} + x_{43} \\
&{}  - 8x_{57}^{2} + 2x_{58} + 4x_{58}^{2} + x_{59} - x_{60} \\
&{}  - 9x_{71} + x_{71}^{2} - x_{72} + 2x_{131}^{2} + 2x_{131} x_{136} \\
&{}  + 2x_{131} x_{161} + 2x_{131} x_{170} + 2x_{133}^{2} - 2x_{134}^{2} + 2x_{135} \\
&{}  - 2x_{137} - x_{138} - x_{142}^{2} + 2x_{142} x_{144} + x_{143}^{2} \\
&{}  - x_{145}^{2} + 2x_{145} x_{147} + x_{146}^{2} - x_{148}^{2} + 2x_{148} x_{150} \\
&{}  + x_{149}^{2} - x_{151}^{2} + 2x_{151} x_{153} + x_{152}^{2} - x_{154}^{2} \\
&{}  + 2x_{154} x_{156} + x_{155}^{2} - x_{157}^{2} + 2x_{157} x_{159} + x_{158}^{2} \\
&{}  + 2x_{160} - 2x_{162} - x_{163}^{2} + 2x_{163} x_{165} + x_{164}^{2} \\
&{}  - x_{166}^{2} + 2x_{166} x_{168} + x_{167}^{2} + 2x_{169} - 2x_{171} \\
&{}  - x_{172}^{2} + 2x_{172} x_{174} + x_{173}^{2} - x_{175}^{2} + 2x_{175} x_{177} \\
&{}  + x_{176}^{2} - 4x_{178}^{2} + x_{179}^{2} + x_{180}^{2} + 2x_{181}^{2} \\
&{}  + 2x_{182}^{2} - x_{189}^{2} - 2x_{190}^{2} + x_{207}^{2} - x_{208}^{2}
\end{align*}
% END RESIDUAL 13
% BEGIN RESIDUAL 14
\begin{align*}
R_{14} ={}& -2 + x_{0} - x_{1} + x_{6} x_{7} + 2x_{8} \\
&{}  - 9x_{8}^{2} + 19x_{8} x_{9} + 5x_{8} x_{10} - 4x_{8} x_{23} - 8x_{8} x_{24} \\
&{}  - 32x_{8} x_{37} - 13x_{9} + 9x_{9}^{2} + 5x_{9} x_{10} - 8x_{9} x_{23} \\
&{}  + 4x_{9} x_{24} + 12x_{10} - 4x_{10}^{2} + 2x_{10} x_{11} - 4x_{10} x_{19} \\
&{}  - 8x_{10} x_{23} - 20x_{10} x_{27} - 24x_{10} x_{37} + x_{11}^{2} - 8x_{19} x_{23} \\
&{}  + 9x_{23} - 8x_{23}^{2} + 16x_{23} x_{24} - x_{24} - 8x_{24}^{2} \\
&{}  - 16x_{24} x_{27} - 8x_{24} x_{37} - 8x_{25}^{2} + 2x_{26} - 2x_{27} \\
&{}  - 4x_{27}^{2} - x_{29} - 4x_{29}^{2} + x_{30} + x_{31} \\
&{}  - 36x_{37} - 4x_{37}^{2} + x_{38} - x_{39} + x_{42} \\
&{}  - 2x_{43} - x_{59} + 7x_{71} - 4x_{178}^{2}
\end{align*}
% END RESIDUAL 14
% BEGIN RESIDUAL 15
\begin{align*}
R_{15} ={}& 2 + x_{0} + 2x_{1} - x_{2} - x_{4} x_{5} \\
&{}  - x_{6} x_{7} + 11x_{8} + 4x_{8}^{2} + 16x_{8} x_{9} - 2x_{8} x_{10} \\
&{}  + 8x_{8} x_{23} - 8x_{8} x_{24} - 4x_{8} x_{39} - 16x_{9} - 9x_{9}^{2} \\
&{}  + 6x_{9} x_{10} - 8x_{9} x_{24} - 32x_{9} x_{36} + 12x_{10} - 7x_{10}^{2} \\
&{}  + 2x_{10} x_{11} + 12x_{10} x_{19} - 16x_{10} x_{23} + 8x_{10} x_{24} - 4x_{10} x_{27} \\
&{}  - 24x_{10} x_{36} - 8x_{10} x_{39} + x_{11}^{2} - 4x_{19} - 8x_{19} x_{23} \\
&{}  + 16x_{19} x_{24} + 8x_{19} x_{27} + x_{20} + 2x_{23} - 4x_{23}^{2} \\
&{}  - 8x_{23} x_{36} - 10x_{24} - 13x_{27} - x_{28} + x_{31} \\
&{}  + 4x_{36} - 4x_{36}^{2} + x_{37} - x_{39} - 4x_{39}^{2} \\
&{}  + x_{40} + x_{42} - x_{43} - x_{59} + 6x_{71} \\
&{}  + x_{131}^{2} + x_{131} x_{136} + x_{131} x_{161} + x_{131} x_{170} + x_{133}^{2} \\
&{}  - x_{134}^{2} + x_{135} - x_{137} + x_{160} - x_{162} \\
&{}  + x_{169} - x_{171} - 4x_{178}^{2} + x_{179}^{2} + x_{180}^{2} \\
&{}  + x_{181}^{2} + x_{182}^{2} - x_{189}^{2} - x_{190}^{2}
\end{align*}
% END RESIDUAL 15
% BEGIN RESIDUAL 16
\begin{align*}
R_{16} ={}& -1 - x_{0} - 17x_{8} + 6x_{8}^{2} + 27x_{8} x_{9} \\
&{}  - 3x_{8} x_{10} - 8x_{8} x_{19} - 60x_{8} x_{23} - 16x_{8} x_{24} - 32x_{8} x_{31} \\
&{}  - 4x_{8} x_{39} + 6x_{9} - 15x_{9}^{2} + 7x_{9} x_{10} + 32x_{9} x_{19} \\
&{}  + 8x_{9} x_{23} + 17x_{10} - 8x_{10}^{2} + 2x_{10} x_{11} - 32x_{10} x_{19} \\
&{}  - 24x_{10} x_{23} + 16x_{10} x_{24} - 20x_{10} x_{27} - 40x_{10} x_{31} - 8x_{10} x_{39} \\
&{}  + x_{11}^{2} + 12x_{19} - 8x_{19} x_{23} - 8x_{19} x_{24} - 4x_{20} \\
&{}  - 5x_{23} - 16x_{23}^{2} + 8x_{23} x_{24} - 16x_{23} x_{31} - 8x_{24} \\
&{}  - 4x_{24}^{2} - 16x_{24} x_{27} - 8x_{24} x_{31} + x_{25} - 4x_{27} \\
&{}  - 4x_{27}^{2} - 2x_{29} - 4x_{29}^{2} + x_{30} - 43x_{31} \\
&{}  - 4x_{31}^{2} + x_{32} + x_{36} - x_{39} - 4x_{39}^{2} \\
&{}  + x_{40} + x_{42} + x_{43} - 2x_{59} + 13x_{71} \\
&{}  + 4x_{178}^{2}
\end{align*}
% END RESIDUAL 16
% BEGIN RESIDUAL 17
\begin{align*}
R_{17} ={}& -4x_{19}^{2} + x_{20}
\end{align*}
% END RESIDUAL 17
% BEGIN RESIDUAL 18
\begin{align*}
R_{18} ={}& x_{8} - x_{8}^{2} + 8x_{8} x_{19} - 4x_{8} x_{20} + 16x_{19} x_{20} \\
&{}  - 4x_{20} - 4x_{20}^{2} + x_{33}
\end{align*}
% END RESIDUAL 18
% BEGIN RESIDUAL 19
\begin{align*}
R_{19} ={}& 1 + 2x_{33} - 4x_{33}^{2} + x_{34}
\end{align*}
% END RESIDUAL 19
% BEGIN RESIDUAL 20
\begin{align*}
R_{20} ={}& -x_{9}^{2} - 4x_{9} x_{34} - 4x_{34}^{2} + x_{35}
\end{align*}
% END RESIDUAL 20
% BEGIN RESIDUAL 21
\begin{align*}
R_{21} ={}& 2 + 2x_{1} - x_{2} - x_{4} x_{5} - x_{6} x_{7} \\
&{}  + 4x_{8} + 6x_{8}^{2} + 5x_{8} x_{9} - x_{8} x_{10} + 16x_{8} x_{19} \\
&{}  - 4x_{9} + 8x_{9}^{2} - 5x_{9} x_{10} + 32x_{9} x_{19} - 11x_{10} \\
&{}  - 3x_{10}^{2} + 2x_{10} x_{11} + 4x_{10} x_{19} + x_{11}^{2} + 25x_{19} \\
&{}  - 3x_{20} - 2x_{23} + x_{36} - x_{39} - x_{59} \\
&{}  + 7x_{71} + x_{131}^{2} + x_{131} x_{136} + x_{133}^{2} - x_{134}^{2} \\
&{}  + x_{135} - x_{137}
\end{align*}
% END RESIDUAL 21
% BEGIN RESIDUAL 22
\begin{align*}
R_{22} ={}& 1 + x_{0} - 3x_{8} - 2x_{8}^{2} - 5x_{8} x_{9} \\
&{}  + x_{8} x_{10} + 11x_{9} - 15x_{9}^{2} - 7x_{9} x_{10} - 8x_{9} x_{23} \\
&{}  - 32x_{9} x_{36} + 16x_{9} x_{43} + 4x_{9} x_{47} + 6x_{10} + 2x_{10}^{2} \\
&{}  - 16x_{10} x_{23} - 24x_{10} x_{36} + 4x_{10} x_{43} + 8x_{10} x_{47} + x_{19} \\
&{}  - 27x_{23} + 8x_{23}^{2} - 8x_{23} x_{36} + 16x_{23} x_{43} + x_{24} \\
&{}  + 4x_{36} - 4x_{36}^{2} + x_{37} + x_{43} + 4x_{43}^{2} \\
&{}  - x_{44} + 4x_{47}^{2} - x_{48} + 3x_{71} - 4x_{178}^{2}
\end{align*}
% END RESIDUAL 22
% BEGIN RESIDUAL 23
\begin{align*}
R_{23} ={}& -4 - x_{0} - 3x_{1} + x_{2} + 2x_{2} x_{208} \\
&{}  + x_{3}^{2} - 2x_{3} x_{131} + x_{4} x_{5} + 2x_{6} x_{7} - 11x_{8} \\
&{}  - 7x_{8}^{2} - 2x_{8} x_{9} - 4x_{8} x_{10} + 24x_{8} x_{19} - 8x_{8} x_{23} \\
&{}  + 17x_{9} + 8x_{9}^{2} + 8x_{9} x_{10} + 24x_{9} x_{23} - 4x_{9} x_{24} \\
&{}  - 16x_{9} x_{43} + 4x_{10}^{2} - 2x_{10} x_{11} + 32x_{10} x_{19} + 8x_{10} x_{23} \\
&{}  - 16x_{10} x_{24} - 4x_{10} x_{43} + 2x_{11} + 21x_{19} + 8x_{19} x_{24} \\
&{}  + x_{20} + 16x_{23} - 8x_{23} x_{24} - 16x_{23} x_{43} - 6x_{24} \\
&{}  - 4x_{24}^{2} - x_{25} - 8x_{25}^{2} + 2x_{26} + x_{27} \\
&{}  + x_{29} + 4x_{29}^{2} - x_{30} + x_{39} - 7x_{43} \\
&{}  - 4x_{43}^{2} + x_{44} - x_{71}^{2} + x_{72} - x_{131}^{2} \\
&{}  - x_{131} x_{136} - x_{131} x_{161} - 2x_{131} x_{170} - x_{133}^{2} + x_{134}^{2} \\
&{}  - x_{135} + x_{137} + x_{138} + x_{142}^{2} - 2x_{142} x_{144} \\
&{}  - x_{143}^{2} + x_{145}^{2} - 2x_{145} x_{147} - x_{146}^{2} + x_{148}^{2} \\
&{}  - 2x_{148} x_{150} - x_{149}^{2} + x_{151}^{2} - 2x_{151} x_{153} - x_{152}^{2} \\
&{}  + x_{154}^{2} - 2x_{154} x_{156} - x_{155}^{2} + x_{157}^{2} - 2x_{157} x_{159} \\
&{}  - x_{158}^{2} - x_{160} + x_{162} + x_{163}^{2} - 2x_{163} x_{165} \\
&{}  - x_{164}^{2} + x_{166}^{2} - 2x_{166} x_{168} - x_{167}^{2} - 2x_{169} \\
&{}  + 2x_{171} + x_{172}^{2} - 2x_{172} x_{174} - x_{173}^{2} + x_{175}^{2} \\
&{}  - 2x_{175} x_{177} - x_{176}^{2} + 4x_{178}^{2} - x_{179}^{2} - x_{180}^{2} \\
&{}  - x_{181}^{2} - x_{182}^{2} + x_{189}^{2} + x_{190}^{2} - x_{207}^{2} \\
&{}  + x_{208}^{2}
\end{align*}
% END RESIDUAL 23
% BEGIN RESIDUAL 24
\begin{align*}
R_{24} ={}& -x_{10}^{2} - 2x_{10} x_{11} - x_{11}^{2} + x_{39}
\end{align*}
% END RESIDUAL 24
% BEGIN RESIDUAL 25
\begin{align*}
R_{25} ={}& -1 - x_{1} + x_{6} x_{7} - 2x_{8} + x_{8} x_{9} \\
&{}  + x_{8} x_{10} - 8x_{8} x_{24} - 4x_{8} x_{39} + x_{9}^{2} + x_{9} x_{10} \\
&{}  - 8x_{9} x_{24} + x_{10}^{2} - 8x_{10} x_{39} - x_{23} - 8x_{23} x_{24} \\
&{}  + 5x_{24} - 4x_{24}^{2} + 4x_{25}^{2} - x_{26} - 4x_{39}^{2} \\
&{}  + x_{40} + x_{43} - x_{71}
\end{align*}
% END RESIDUAL 25
% BEGIN RESIDUAL 26
\begin{align*}
R_{26} ={}& -2x_{8} - 2x_{8}^{2} + 3x_{8} x_{9} - 11x_{8} x_{10} + 8x_{8} x_{23} \\
&{}  - 8x_{8} x_{24} - 28x_{8} x_{39} - 16x_{8} x_{40} + 4x_{8} x_{52} + 4x_{9} \\
&{}  - x_{9}^{2} + 3x_{9} x_{10} + 24x_{9} x_{24} - 5x_{10} - 5x_{10}^{2} \\
&{}  - 16x_{10} x_{23} - 8x_{10} x_{24} + 20x_{10} x_{27} - 24x_{10} x_{39} - 16x_{10} x_{40} \\
&{}  + 8x_{10} x_{52} - x_{19} + 2x_{23} - 4x_{23}^{2} + 16x_{23} x_{24} \\
&{}  - 27x_{24} + 8x_{24}^{2} + 16x_{24} x_{27} - 32x_{24} x_{39} - 16x_{24} x_{40} \\
&{}  - 8x_{25}^{2} + 2x_{26} + x_{27} + 4x_{27}^{2} - 4x_{39} \\
&{}  - 12x_{39}^{2} - 16x_{39} x_{40} - 3x_{40} - 4x_{40}^{2} + x_{41} \\
&{}  - x_{42} - x_{43} + 4x_{52}^{2} - x_{53} + 7x_{71}
\end{align*}
% END RESIDUAL 26
% BEGIN RESIDUAL 27
\begin{align*}
R_{27} ={}& 2 + x_{1} - x_{6} x_{7} + 2x_{8}^{2} + 2x_{8} x_{9} \\
&{}  + 2x_{8} x_{10} + x_{9}^{2} + x_{10} - 3x_{10}^{2} + 2x_{10} x_{11} \\
&{}  + 16x_{10} x_{23} - 24x_{10} x_{24} - 20x_{10} x_{27} + x_{11}^{2} - x_{23} \\
&{}  + 4x_{23}^{2} + 8x_{23} x_{24} - 2x_{24} - 12x_{24}^{2} - 16x_{24} x_{27} \\
&{}  - x_{27} - 4x_{27}^{2} - x_{39} + x_{42} - 2x_{71} \\
&{}  + x_{131}^{2} + x_{133}^{2} - x_{134}^{2}
\end{align*}
% END RESIDUAL 27
% BEGIN RESIDUAL 28
\begin{align*}
R_{28} ={}& -1 - 3x_{8} - x_{8}^{2} - x_{8} x_{9} - x_{8} x_{10} \\
&{}  + 4x_{9} - 3x_{9}^{2} + 3x_{9} x_{10} - 8x_{9} x_{23} + 3x_{10}^{2} \\
&{}  + x_{19} + 6x_{23} - 4x_{23}^{2} + x_{43} - x_{71}
\end{align*}
% END RESIDUAL 28
% BEGIN RESIDUAL 29
\begin{align*}
R_{29} ={}& 2 + 2x_{1} - x_{2} - x_{4} x_{5} - x_{6} x_{7} \\
&{}  + 6x_{8} + 2x_{8}^{2} + 2x_{8} x_{9} + 2x_{8} x_{10} - 8x_{9} \\
&{}  - 6x_{9}^{2} - 6x_{9} x_{10} - 16x_{9} x_{23} - 16x_{9} x_{43} - 3x_{10}^{2} \\
&{}  - 8x_{10} x_{23} - 4x_{10} x_{43} - 2x_{19} - 16x_{23} - 8x_{23}^{2} \\
&{}  - 16x_{23} x_{43} - 2x_{43} - 4x_{43}^{2} + x_{44} + 2x_{71} \\
&{}  + x_{131}^{2} + x_{131} x_{136} + x_{131} x_{161} + x_{131} x_{170} + x_{133}^{2} \\
&{}  - x_{134}^{2} + x_{135} - x_{137} + x_{160} - x_{162} \\
&{}  + x_{169} - x_{171} + x_{179}^{2} + x_{180}^{2} - x_{189}^{2}
\end{align*}
% END RESIDUAL 29
% BEGIN RESIDUAL 30
\begin{align*}
R_{30} ={}& 1 + x_{1} - x_{6} x_{7} + x_{8}^{2} + 2x_{8} x_{10} \\
&{}  + x_{9} - 2x_{9} x_{10} - 4x_{9} x_{44} - 2x_{10}^{2} - 8x_{10} x_{44} \\
&{}  - x_{23} - x_{24} - 4x_{44}^{2} + x_{45}
\end{align*}
% END RESIDUAL 30
% BEGIN RESIDUAL 31
\begin{align*}
R_{31} ={}& -1 - 4x_{45} - 4x_{45}^{2} + x_{46}
\end{align*}
% END RESIDUAL 31
% BEGIN RESIDUAL 32
\begin{align*}
R_{32} ={}& 1 - x_{10}^{2} - 8x_{10} x_{45} - 4x_{10} x_{46} + 8x_{45} \\
&{}  - 8x_{45}^{2} - 16x_{45} x_{46} - 2x_{46} - 4x_{46}^{2} + x_{47}
\end{align*}
% END RESIDUAL 32
% BEGIN RESIDUAL 33
\begin{align*}
R_{33} ={}& 2 + x_{1} - x_{6} x_{7} + x_{8}^{2} + 2x_{8} x_{10} \\
&{}  + x_{9} - 2x_{9} x_{10} - 4x_{9} x_{47} - 2x_{10}^{2} - 8x_{10} x_{47} \\
&{}  - x_{23} - x_{24} - 4x_{47}^{2} + x_{48}
\end{align*}
% END RESIDUAL 33
% BEGIN RESIDUAL 34
\begin{align*}
R_{34} ={}& 1 - 4x_{42}^{2} - 8x_{42} x_{48} - 4x_{48}^{2} + x_{49}
\end{align*}
% END RESIDUAL 34
% BEGIN RESIDUAL 35
\begin{align*}
R_{35} ={}& x_{8} - x_{8}^{2} - 8x_{8} x_{42} - 4x_{8} x_{49} - 12x_{42}^{2} \\
&{}  + 8x_{42} x_{48} - 16x_{42} x_{49} + 4x_{48}^{2} - x_{49} - 4x_{49}^{2} \\
&{}  + x_{50}
\end{align*}
% END RESIDUAL 35
% BEGIN RESIDUAL 36
\begin{align*}
R_{36} ={}& -1 - 4x_{50} - 4x_{50}^{2} + x_{51}
\end{align*}
% END RESIDUAL 36
% BEGIN RESIDUAL 37
\begin{align*}
R_{37} ={}& 1 - x_{10}^{2} - 8x_{10} x_{50} - 4x_{10} x_{51} + 8x_{50} \\
&{}  - 8x_{50}^{2} - 16x_{50} x_{51} - 2x_{51} - 4x_{51}^{2} + x_{52}
\end{align*}
% END RESIDUAL 37
% BEGIN RESIDUAL 38
\begin{align*}
R_{38} ={}& x_{8} + x_{8}^{2} - 4x_{8} x_{52} - 2x_{10}^{2} - 8x_{10} x_{52} \\
&{}  - 2x_{24} - 4x_{52}^{2} + x_{53}
\end{align*}
% END RESIDUAL 38
% BEGIN RESIDUAL 39
\begin{align*}
R_{39} ={}& -2 - 2x_{1} + x_{2} + x_{4} x_{5} + x_{6} x_{7} \\
&{}  - x_{8}^{2} - 2x_{8} x_{10} + x_{9}^{2} + 2x_{9} x_{10} - x_{23} \\
&{}  + x_{24} - 12x_{53} - 4x_{53}^{2} + x_{54} - x_{131}^{2} \\
&{}  - x_{131} x_{136} - x_{131} x_{161} - x_{133}^{2} + x_{134}^{2} - x_{135} \\
&{}  + x_{137} - x_{160} + x_{162}
\end{align*}
% END RESIDUAL 39
% BEGIN RESIDUAL 40
\begin{align*}
R_{40} ={}& -4x_{9} - x_{9}^{2} - 2x_{9} x_{10} - 16x_{9} x_{53} - 8x_{9} x_{54} \\
&{}  - x_{10}^{2} - 3x_{23} - 16x_{23} x_{53} - 8x_{23} x_{54} - x_{43} \\
&{}  + 8x_{53} - 16x_{53}^{2} - 16x_{53} x_{54} + 4x_{54} - 4x_{54}^{2} \\
&{}  + x_{55}
\end{align*}
% END RESIDUAL 40
% BEGIN RESIDUAL 41
\begin{align*}
R_{41} ={}& -2 - 2x_{1} + x_{2} + x_{4} x_{5} + x_{6} x_{7} \\
&{}  + 6x_{8} + 2x_{8}^{2} + 4x_{8} x_{9} + x_{9} - 5x_{9}^{2} \\
&{}  + 4x_{9} x_{10} + 16x_{9} x_{43} - 20x_{9} x_{55} - x_{10}^{2} + 8x_{10} x_{23} \\
&{}  + 4x_{10} x_{43} - 2x_{19} - 6x_{23} + 4x_{23}^{2} + 16x_{23} x_{43} \\
&{}  - 16x_{23} x_{55} - x_{43} + 4x_{43}^{2} - x_{44} - 12x_{48} \\
&{}  - 4x_{48}^{2} - 4x_{55}^{2} + x_{56}
\end{align*}
% END RESIDUAL 41
% BEGIN RESIDUAL 42
\begin{align*}
R_{42} ={}& 3 + 2x_{1} - x_{2} - x_{4} x_{5} - x_{6} x_{7} \\
&{}  - 3x_{8} - 2x_{8} x_{9} + 2x_{8} x_{10} - 4x_{9} - 2x_{9}^{2} \\
&{}  - 2x_{9} x_{10} + x_{19} + x_{23} - x_{24} + 24x_{48} \\
&{}  - 8x_{48}^{2} - 16x_{48} x_{56} - 2x_{56} - 4x_{56}^{2} + x_{57} \\
&{}  + x_{131}^{2} + x_{131} x_{136} + x_{131} x_{161} + x_{131} x_{170} + x_{133}^{2} \\
&{}  - x_{134}^{2} + x_{135} - x_{137} - x_{154}^{2} + 2x_{154} x_{156} \\
&{}  + x_{155}^{2} - x_{157}^{2} + 2x_{157} x_{159} + x_{158}^{2} + x_{160} \\
&{}  - x_{162} - x_{163}^{2} + 2x_{163} x_{165} + x_{164}^{2} + x_{169} \\
&{}  - x_{171} + x_{179}^{2} + x_{180}^{2} + x_{181}^{2} + x_{182}^{2} \\
&{}  - x_{189}^{2} - x_{190}^{2}
\end{align*}
% END RESIDUAL 42
% BEGIN RESIDUAL 43
\begin{align*}
R_{43} ={}& -2 - x_{1} + x_{6} x_{7} + x_{8} - x_{8}^{2} \\
&{}  - x_{8} x_{9} - x_{8} x_{10} + 4x_{9} + x_{9} x_{10} + x_{10} \\
&{}  + x_{10}^{2} - 4x_{10} x_{57} - x_{19} - x_{23} - 4x_{57}^{2} \\
&{}  + x_{58} + 12x_{59} - 4x_{59}^{2} - x_{71}
\end{align*}
% END RESIDUAL 43
% BEGIN RESIDUAL 44
\begin{align*}
R_{44} ={}& 2 + x_{0} + x_{1} - x_{6} x_{7} + 15x_{8} \\
&{}  + 6x_{8}^{2} - 6x_{8} x_{9} - 2x_{8} x_{10} + 16x_{8} x_{19} + 8x_{8} x_{23} \\
&{}  - 38x_{9}^{2} + 8x_{9} x_{10} + 32x_{9} x_{19} + 8x_{9} x_{23} + 4x_{9} x_{34} \\
&{}  + 11x_{10} + 5x_{10}^{2} - 2x_{10} x_{11} - 4x_{10} x_{19} - x_{11}^{2} \\
&{}  - 2x_{19} - x_{20} - 8x_{23} + 4x_{23}^{2} + x_{24} \\
&{}  + 4x_{25}^{2} - x_{26} + 4x_{34}^{2} - x_{35} + x_{39} \\
&{}  + x_{43} + x_{59} - 8x_{71} + x_{131}^{2} + x_{133}^{2} \\
&{}  - x_{134}^{2} - 4x_{178}^{2}
\end{align*}
% END RESIDUAL 44
% BEGIN RESIDUAL 45
\begin{align*}
R_{45} ={}& -1 - x_{1} + x_{6} x_{7} - 2x_{10} + 2x_{10}^{2} \\
&{}  + 8x_{10} x_{57} + 8x_{57}^{2} - 2x_{58} - 4x_{58}^{2} + x_{60}
\end{align*}
% END RESIDUAL 45
% BEGIN RESIDUAL 46
\begin{align*}
R_{46} ={}& -1 - 4x_{60} - 4x_{60}^{2} + x_{61}
\end{align*}
% END RESIDUAL 46
% BEGIN RESIDUAL 47
\begin{align*}
R_{47} ={}& -x_{2} x_{13} + x_{62}
\end{align*}
% END RESIDUAL 47
% BEGIN RESIDUAL 48
\begin{align*}
R_{48} ={}& x_{12} - x_{13} - 2x_{22} x_{62} - x_{63}
\end{align*}
% END RESIDUAL 48
% BEGIN RESIDUAL 49
\begin{align*}
R_{49} ={}& x_{12} - x_{13} - 2x_{26} x_{62} - x_{64}
\end{align*}
% END RESIDUAL 49
% BEGIN RESIDUAL 50
\begin{align*}
R_{50} ={}& -x_{2} x_{13} + x_{12} - x_{13} - 2x_{28} x_{62} - x_{62} \\
&{}  - x_{65}
\end{align*}
% END RESIDUAL 50
% BEGIN RESIDUAL 51
\begin{align*}
R_{51} ={}& x_{12} - x_{13} - 2x_{30} x_{62} - x_{66}
\end{align*}
% END RESIDUAL 51
% BEGIN RESIDUAL 52
\begin{align*}
R_{52} ={}& x_{2} x_{13} + x_{12} - x_{13} - 2x_{32} x_{62} + x_{62} \\
&{}  - x_{67}
\end{align*}
% END RESIDUAL 52
% BEGIN RESIDUAL 53
\begin{align*}
R_{53} ={}& x_{12} - x_{13} - 4x_{34} x_{62} - 2x_{35} x_{62} - x_{68}
\end{align*}
% END RESIDUAL 53
% BEGIN RESIDUAL 54
\begin{align*}
R_{54} ={}& x_{12} - x_{13} - 2x_{38} x_{62} - x_{69}
\end{align*}
% END RESIDUAL 54
% BEGIN RESIDUAL 55
\begin{align*}
R_{55} ={}& -x_{2} x_{13} + 2x_{28} x_{62} - 2x_{41} x_{62} - x_{62} + x_{65} \\
&{}  - x_{70}
\end{align*}
% END RESIDUAL 55
% BEGIN RESIDUAL 56
\begin{align*}
R_{56} ={}& -x_{8} x_{9} + x_{8} x_{10} + x_{9} x_{10} - x_{10}^{2} + x_{71}
\end{align*}
% END RESIDUAL 56
% BEGIN RESIDUAL 57
\begin{align*}
R_{57} ={}& 1 + 2x_{11} + x_{11}^{2} - x_{71}^{2} + x_{72}
\end{align*}
% END RESIDUAL 57
% BEGIN RESIDUAL 58
\begin{align*}
R_{58} ={}& -x_{12}^{2} + 2x_{12} x_{13} + 2x_{12} x_{74} - x_{13}^{2} - 2x_{13} x_{74} \\
&{}  - x_{72}^{2} + x_{73} - x_{74}^{2}
\end{align*}
% END RESIDUAL 58
% BEGIN RESIDUAL 59
\begin{align*}
R_{59} ={}& -4x_{60} x_{62} - 2x_{61} x_{62} + x_{74}
\end{align*}
% END RESIDUAL 59
% BEGIN RESIDUAL 60
\begin{align*}
R_{60} ={}& -x_{9} x_{62} - x_{13} + x_{75} x_{76} + x_{77}
\end{align*}
% END RESIDUAL 60
% BEGIN RESIDUAL 61
\begin{align*}
R_{61} ={}& -2 + 4x_{2} + 4x_{2} x_{8} + 8x_{2} x_{9} - 2x_{2} x_{19} \\
&{}  + x_{3} x_{18} - 2x_{4} x_{5} + 2x_{6} x_{7} - 2x_{16} + x_{75} \\
&{}  + 2x_{209} x_{211} + x_{211}
\end{align*}
% END RESIDUAL 61
% BEGIN RESIDUAL 62
\begin{align*}
R_{62} ={}& -x_{8} x_{201} - x_{12} + x_{76}
\end{align*}
% END RESIDUAL 62
% BEGIN RESIDUAL 63
\begin{align*}
R_{63} ={}& x_{2} x_{13} + x_{9} x_{201} + x_{12} - x_{13} - 2x_{19} x_{62} \\
&{}  + x_{62} + x_{78}
\end{align*}
% END RESIDUAL 63
% BEGIN RESIDUAL 64
\begin{align*}
R_{64} ={}& -x_{12} x_{63} + x_{13} x_{63} + x_{79}
\end{align*}
% END RESIDUAL 64
% BEGIN RESIDUAL 65
\begin{align*}
R_{65} ={}& -x_{64} x_{79} + x_{80}
\end{align*}
% END RESIDUAL 65
% BEGIN RESIDUAL 66
\begin{align*}
R_{66} ={}& -x_{65} x_{80} + x_{81}
\end{align*}
% END RESIDUAL 66
% BEGIN RESIDUAL 67
\begin{align*}
R_{67} ={}& -x_{66} x_{81} + x_{82}
\end{align*}
% END RESIDUAL 67
% BEGIN RESIDUAL 68
\begin{align*}
R_{68} ={}& -x_{67} x_{82} + x_{83}
\end{align*}
% END RESIDUAL 68
% BEGIN RESIDUAL 69
\begin{align*}
R_{69} ={}& -x_{68} x_{83} + x_{84}
\end{align*}
% END RESIDUAL 69
% BEGIN RESIDUAL 70
\begin{align*}
R_{70} ={}& -x_{69} x_{84} + x_{85}
\end{align*}
% END RESIDUAL 70
% BEGIN RESIDUAL 71
\begin{align*}
R_{71} ={}& -x_{70} x_{85} + x_{86}
\end{align*}
% END RESIDUAL 71
% BEGIN RESIDUAL 72
\begin{align*}
R_{72} ={}& -x_{77} x_{86} + x_{87}
\end{align*}
% END RESIDUAL 72
% BEGIN RESIDUAL 73
\begin{align*}
R_{73} ={}& -x_{78} x_{87} + x_{88}
\end{align*}
% END RESIDUAL 73
% BEGIN RESIDUAL 74
\begin{align*}
R_{74} ={}& -x_{88}^{2} + x_{89}
\end{align*}
% END RESIDUAL 74
% BEGIN RESIDUAL 75
\begin{align*}
R_{75} ={}& -1 - x_{2} + x_{3} x_{18} - x_{4} - 2x_{4} x_{16} \\
&{}  + x_{90}
\end{align*}
% END RESIDUAL 75
% BEGIN RESIDUAL 76
\begin{align*}
R_{76} ={}& -x_{14} x_{90} + x_{91}
\end{align*}
% END RESIDUAL 76
% BEGIN RESIDUAL 77
\begin{align*}
R_{77} ={}& -2x_{1}^{2} + 2x_{1} x_{12} + 2x_{1} x_{13} + 2x_{1} x_{91} + 2x_{1} x_{92} \\
&{}  - x_{12}^{2} - 2x_{12} x_{91} - x_{13}^{2} - 2x_{13} x_{92} + x_{17} x_{93} \\
&{}  - x_{73} x_{89} - x_{91}^{2} - x_{92}^{2}
\end{align*}
% END RESIDUAL 77
% BEGIN RESIDUAL 78
\begin{align*}
R_{78} ={}& -x_{4} x_{15} + x_{15} - x_{15} x_{90} + x_{92}
\end{align*}
% END RESIDUAL 78
% BEGIN RESIDUAL 79
\begin{align*}
R_{79} ={}& -x_{94} x_{97} + x_{97} + x_{99}
\end{align*}
% END RESIDUAL 79
% BEGIN RESIDUAL 80
\begin{align*}
R_{80} ={}& -x_{3} x_{120} + x_{95} x_{193} - x_{98} - x_{99} x_{120}
\end{align*}
% END RESIDUAL 80
% BEGIN RESIDUAL 81
\begin{align*}
R_{81} ={}& x_{100} x_{193} - x_{101} x_{123} - x_{102}
\end{align*}
% END RESIDUAL 81
% BEGIN RESIDUAL 82
\begin{align*}
R_{82} ={}& x_{103} x_{193} - x_{104} x_{123} - x_{105}
\end{align*}
% END RESIDUAL 82
% BEGIN RESIDUAL 83
\begin{align*}
R_{83} ={}& x_{106} x_{193} - x_{107} x_{123} - x_{108}
\end{align*}
% END RESIDUAL 83
% BEGIN RESIDUAL 84
\begin{align*}
R_{84} ={}& x_{109} x_{193} - x_{110} x_{123} - x_{111}
\end{align*}
% END RESIDUAL 84
% BEGIN RESIDUAL 85
\begin{align*}
R_{85} ={}& x_{112} x_{194} - x_{114} x_{123} - x_{115}
\end{align*}
% END RESIDUAL 85
% BEGIN RESIDUAL 86
\begin{align*}
R_{86} ={}& x_{116} x_{194} - x_{117} x_{123} - x_{118}
\end{align*}
% END RESIDUAL 86
% BEGIN RESIDUAL 87
\begin{align*}
R_{87} ={}& -x_{17} x_{120} + x_{123}
\end{align*}
% END RESIDUAL 87
% BEGIN RESIDUAL 88
\begin{align*}
R_{88} ={}& x_{119} - x_{121} x_{123} - x_{122}
\end{align*}
% END RESIDUAL 88
% BEGIN RESIDUAL 89
\begin{align*}
R_{89} ={}& -x_{123} x_{125} + x_{124} - x_{126}
\end{align*}
% END RESIDUAL 89
% BEGIN RESIDUAL 90
\begin{align*}
R_{90} ={}& -x_{123} x_{128} + x_{127} x_{193} - x_{129}
\end{align*}
% END RESIDUAL 90
% BEGIN RESIDUAL 91
\begin{align*}
R_{91} ={}& x_{130} - 2x_{130} x_{132} + x_{133}
\end{align*}
% END RESIDUAL 91
% BEGIN RESIDUAL 92
\begin{align*}
R_{92} ={}& -x_{130} x_{133} - x_{133} + x_{134}
\end{align*}
% END RESIDUAL 92
% BEGIN RESIDUAL 93
\begin{align*}
R_{93} ={}& -1 + x_{131}^{2} + x_{133}^{2} - x_{134}^{2}
\end{align*}
% END RESIDUAL 93
% BEGIN RESIDUAL 94
\begin{align*}
R_{94} ={}& 1 - x_{131} x_{136} - x_{135} + x_{137}
\end{align*}
% END RESIDUAL 94
% BEGIN RESIDUAL 95
\begin{align*}
R_{95} ={}& -x_{131} x_{136} + x_{135} - x_{135}^{2} + 2x_{135} x_{137} + x_{136}^{2} \\
&{}  + x_{137}
\end{align*}
% END RESIDUAL 95
% BEGIN RESIDUAL 96
\begin{align*}
R_{96} ={}& 1 + x_{3}^{2} - 2x_{3} x_{131} + x_{138}
\end{align*}
% END RESIDUAL 96
% BEGIN RESIDUAL 97
\begin{align*}
R_{97} ={}& -x_{3} x_{170} + x_{17} + x_{138} x_{139} - x_{169}
\end{align*}
% END RESIDUAL 97
% BEGIN RESIDUAL 98
\begin{align*}
R_{98} ={}& 2 - 2x_{2} - 4x_{2}^{2} - 16x_{2} x_{16} + 4x_{2} x_{131} \\
&{}  - 5x_{3} x_{18} - x_{4} + x_{4} x_{5} - 2x_{4} x_{16} - x_{6} x_{7} \\
&{}  - 4x_{16} - 8x_{16}^{2} + 8x_{16} x_{131} + x_{90} + 3x_{131} \\
&{}  + 2x_{136}^{2} - x_{140} + 2x_{143}^{2} + 2x_{146}^{2} + 2x_{149}^{2} \\
&{}  + 2x_{152}^{2} + 2x_{155}^{2} + 2x_{158}^{2} + 2x_{161}^{2} + 2x_{164}^{2} \\
&{}  + 2x_{167}^{2} + 2x_{170}^{2} + 2x_{173}^{2} - x_{176}^{2} - x_{179} x_{186} \\
&{}  + x_{183} - x_{191} + 2x_{195}^{2} + 2x_{196}^{2} + 2x_{197} \\
&{}  + 2x_{197}^{2}
\end{align*}
% END RESIDUAL 98
% BEGIN RESIDUAL 99
\begin{align*}
R_{99} ={}& -1 - 2x_{2} x_{170} - 4x_{16} x_{170} + x_{94} + x_{140} x_{141} \\
&{}  - x_{169} - 2x_{170}
\end{align*}
% END RESIDUAL 99
% BEGIN RESIDUAL 100
\begin{align*}
R_{100} ={}& -x_{131} x_{143} - x_{142} + x_{144}
\end{align*}
% END RESIDUAL 100
% BEGIN RESIDUAL 101
\begin{align*}
R_{101} ={}& -1 - x_{142}^{2} + 2x_{142} x_{144} + x_{143}^{2}
\end{align*}
% END RESIDUAL 101
% BEGIN RESIDUAL 102
\begin{align*}
R_{102} ={}& -x_{131} x_{146} - x_{145} + x_{147}
\end{align*}
% END RESIDUAL 102
% BEGIN RESIDUAL 103
\begin{align*}
R_{103} ={}& -1 - x_{145}^{2} + 2x_{145} x_{147} + x_{146}^{2}
\end{align*}
% END RESIDUAL 103
% BEGIN RESIDUAL 104
\begin{align*}
R_{104} ={}& -x_{131} x_{149} - x_{148} + x_{150}
\end{align*}
% END RESIDUAL 104
% BEGIN RESIDUAL 105
\begin{align*}
R_{105} ={}& -1 - x_{148}^{2} + 2x_{148} x_{150} + x_{149}^{2}
\end{align*}
% END RESIDUAL 105
% BEGIN RESIDUAL 106
\begin{align*}
R_{106} ={}& -x_{131} x_{152} - x_{151} + x_{153}
\end{align*}
% END RESIDUAL 106
% BEGIN RESIDUAL 107
\begin{align*}
R_{107} ={}& -1 - x_{151}^{2} + 2x_{151} x_{153} + x_{152}^{2}
\end{align*}
% END RESIDUAL 107
% BEGIN RESIDUAL 108
\begin{align*}
R_{108} ={}& -x_{131} x_{155} - x_{154} + x_{156}
\end{align*}
% END RESIDUAL 108
% BEGIN RESIDUAL 109
\begin{align*}
R_{109} ={}& -1 - x_{154}^{2} + 2x_{154} x_{156} + x_{155}^{2}
\end{align*}
% END RESIDUAL 109
% BEGIN RESIDUAL 110
\begin{align*}
R_{110} ={}& -x_{131} x_{158} - x_{157} + x_{159}
\end{align*}
% END RESIDUAL 110
% BEGIN RESIDUAL 111
\begin{align*}
R_{111} ={}& -1 - x_{157}^{2} + 2x_{157} x_{159} + x_{158}^{2}
\end{align*}
% END RESIDUAL 111
% BEGIN RESIDUAL 112
\begin{align*}
R_{112} ={}& 1 - x_{131} x_{161} - x_{160} + x_{162}
\end{align*}
% END RESIDUAL 112
% BEGIN RESIDUAL 113
\begin{align*}
R_{113} ={}& -x_{131} x_{161} + x_{160} - x_{160}^{2} + 2x_{160} x_{162} + x_{161}^{2} \\
&{}  + x_{162}
\end{align*}
% END RESIDUAL 113
% BEGIN RESIDUAL 114
\begin{align*}
R_{114} ={}& -x_{131} x_{164} - x_{163} + x_{165}
\end{align*}
% END RESIDUAL 114
% BEGIN RESIDUAL 115
\begin{align*}
R_{115} ={}& -1 - x_{163}^{2} + 2x_{163} x_{165} + x_{164}^{2}
\end{align*}
% END RESIDUAL 115
% BEGIN RESIDUAL 116
\begin{align*}
R_{116} ={}& -x_{131} x_{167} - x_{166} + x_{168}
\end{align*}
% END RESIDUAL 116
% BEGIN RESIDUAL 117
\begin{align*}
R_{117} ={}& -1 - x_{166}^{2} + 2x_{166} x_{168} + x_{167}^{2}
\end{align*}
% END RESIDUAL 117
% BEGIN RESIDUAL 118
\begin{align*}
R_{118} ={}& 1 - x_{131} x_{170} - x_{169} + x_{171}
\end{align*}
% END RESIDUAL 118
% BEGIN RESIDUAL 119
\begin{align*}
R_{119} ={}& -x_{131} x_{170} + x_{169} - x_{169}^{2} + 2x_{169} x_{171} + x_{170}^{2} \\
&{}  + x_{171}
\end{align*}
% END RESIDUAL 119
% BEGIN RESIDUAL 120
\begin{align*}
R_{120} ={}& -x_{131} x_{173} - x_{172} + x_{174}
\end{align*}
% END RESIDUAL 120
% BEGIN RESIDUAL 121
\begin{align*}
R_{121} ={}& -1 - x_{172}^{2} + 2x_{172} x_{174} + x_{173}^{2}
\end{align*}
% END RESIDUAL 121
% BEGIN RESIDUAL 122
\begin{align*}
R_{122} ={}& -x_{131} x_{176} - x_{175} + x_{177}
\end{align*}
% END RESIDUAL 122
% BEGIN RESIDUAL 123
\begin{align*}
R_{123} ={}& -1 - x_{175}^{2} + 2x_{175} x_{177} + x_{176}^{2}
\end{align*}
% END RESIDUAL 123
% BEGIN RESIDUAL 124
\begin{align*}
R_{124} ={}& -1 - 2x_{0} - 2x_{0}^{2} - 2x_{0} x_{2} - x_{1}^{2} \\
&{}  - 2x_{2}^{2} - 4x_{2} x_{131} - x_{3}^{2} - x_{4}^{2} - x_{8}^{2} \\
&{}  + 2x_{9} x_{10} - x_{11}^{2} - x_{12}^{2} - x_{13}^{2} - x_{14}^{2} \\
&{}  - x_{15}^{2} - 4x_{16}^{2} - 8x_{16} x_{131} - x_{17}^{2} - x_{23} \\
&{}  + 2x_{94} - 2x_{94}^{2} - x_{95}^{2} - x_{96}^{2} - x_{100}^{2} \\
&{}  - x_{103}^{2} - x_{106}^{2} - x_{109}^{2} - x_{112}^{2} - x_{113}^{2} \\
&{}  - x_{116}^{2} - x_{119}^{2} - x_{120}^{2} - x_{124}^{2} - x_{127}^{2} \\
&{}  + x_{130} - 3x_{131} + x_{140} + x_{179} x_{186} - x_{183} \\
&{}  - 2x_{210}
\end{align*}
% END RESIDUAL 124
% BEGIN RESIDUAL 125
\begin{align*}
R_{125} ={}& -8 + x_{0} + x_{0}^{2} + 2x_{0} x_{2} - 2x_{0} x_{136} \\
&{}  - 2x_{0} x_{143} - 4x_{1} + x_{1}^{2} - 2x_{1} x_{170} + 3x_{2} \\
&{}  + x_{2}^{2} - 2x_{2} x_{143} + 2x_{2} x_{206} + 2x_{2} x_{208} + x_{3}^{2} \\
&{}  - 2x_{3} x_{18} - 2x_{3} x_{131} + x_{3} x_{170} + x_{4}^{2} + x_{4} x_{5} \\
&{}  + 2x_{4} x_{16} - 2x_{4} x_{161} + 3x_{6} x_{7} + 3x_{8} + x_{8} x_{9} \\
&{}  - 3x_{8} x_{10} - 2x_{8} x_{146} + 4x_{9} + 4x_{9}^{2} + 5x_{9} x_{10} \\
&{}  - 2x_{9} x_{149} - 2x_{10} x_{152} + 2x_{11} + 2x_{11}^{2} - 2x_{11} x_{155} \\
&{}  + x_{12}^{2} - 2x_{12} x_{158} + x_{13}^{2} - 2x_{13} x_{164} + x_{14}^{2} \\
&{}  - 2x_{14} x_{167} + x_{15}^{2} - 2x_{15} x_{173} + 2x_{16} + 4x_{16}^{2} \\
&{}  - 4x_{16} x_{176} - x_{17} - x_{19} - 2x_{23} + 2x_{24} \\
&{}  + x_{71} - x_{71}^{2} + x_{72} - x_{90} - 3x_{131}^{2} \\
&{}  - 3x_{131} x_{136} - x_{131} x_{143} - x_{131} x_{146} - x_{131} x_{149} - x_{131} x_{152} \\
&{}  - x_{131} x_{155} - x_{131} x_{158} - 3x_{131} x_{161} - x_{131} x_{164} - x_{131} x_{167} \\
&{}  - 2x_{131} x_{170} - x_{131} x_{173} - x_{131} x_{176} - x_{131} x_{192} - 3x_{133}^{2} \\
&{}  + 3x_{134}^{2} - 3x_{135} + x_{136}^{2} + 3x_{137} + x_{137}^{2} \\
&{}  + x_{138} - x_{138} x_{139} - x_{142} + x_{142}^{2} - 2x_{142} x_{144} \\
&{}  - x_{144} + x_{144}^{2} - x_{145} + x_{145}^{2} - 2x_{145} x_{147} \\
&{}  - x_{147} + x_{147}^{2} - x_{148} + x_{148}^{2} - 2x_{148} x_{150} \\
&{}  - x_{150} + x_{150}^{2} - x_{151} + x_{151}^{2} - 2x_{151} x_{153} \\
&{}  - x_{153} + x_{153}^{2} - x_{154} + x_{154}^{2} - 2x_{154} x_{156} \\
&{}  - x_{156} + x_{156}^{2} - x_{157} + x_{157}^{2} - 2x_{157} x_{159} \\
&{}  - x_{159} + x_{159}^{2} - 3x_{160} + x_{161}^{2} + 3x_{162} \\
&{}  + x_{162}^{2} - x_{163} + x_{163}^{2} - 2x_{163} x_{165} - x_{165} \\
&{}  + x_{165}^{2} - x_{166} + x_{166}^{2} - 2x_{166} x_{168} - x_{168} \\
&{}  + x_{168}^{2} - 3x_{169} + x_{169}^{2} - 2x_{169} x_{171} + 2x_{171} \\
&{}  + x_{171}^{2} - x_{172} + x_{172}^{2} - 2x_{172} x_{174} - x_{174} \\
&{}  + x_{174}^{2} - x_{175} + x_{175}^{2} - 2x_{175} x_{177} - x_{177} \\
&{}  + x_{177}^{2} - 4x_{178}^{2} - 2x_{179}^{2} - 2x_{180}^{2} - 2x_{181}^{2} \\
&{}  - 2x_{182}^{2} + 2x_{189}^{2} + 2x_{190}^{2} + x_{192} - x_{207}^{2} \\
&{}  + x_{208}^{2} + 2x_{209} x_{211} + x_{210} + x_{211}
\end{align*}
% END RESIDUAL 125
% BEGIN RESIDUAL 126
\begin{align*}
R_{126} ={}& -x_{131} x_{180} + x_{189}
\end{align*}
% END RESIDUAL 126
% BEGIN RESIDUAL 127
\begin{align*}
R_{127} ={}& -1 + x_{179}^{2} + x_{180}^{2} - x_{189}^{2}
\end{align*}
% END RESIDUAL 127
% BEGIN RESIDUAL 128
\begin{align*}
R_{128} ={}& -x_{182} x_{183} + x_{190}
\end{align*}
% END RESIDUAL 128
% BEGIN RESIDUAL 129
\begin{align*}
R_{129} ={}& -1 + x_{181}^{2} + x_{182}^{2} - x_{190}^{2}
\end{align*}
% END RESIDUAL 129
% BEGIN RESIDUAL 130
\begin{align*}
R_{130} ={}& -x_{176}^{2} + x_{191}
\end{align*}
% END RESIDUAL 130
% BEGIN RESIDUAL 131
\begin{align*}
R_{131} ={}& x_{180} - 2x_{185} x_{191} - x_{191}
\end{align*}
% END RESIDUAL 131
% BEGIN RESIDUAL 132
\begin{align*}
R_{132} ={}& -1 - 3x_{180} - 8x_{180} x_{184} + x_{183} - 2x_{185} x_{191} \\
&{}  - x_{191}
\end{align*}
% END RESIDUAL 132
% BEGIN RESIDUAL 133
\begin{align*}
R_{133} ={}& -x_{131} - x_{179} x_{186} + x_{183}
\end{align*}
% END RESIDUAL 133
% BEGIN RESIDUAL 134
\begin{align*}
R_{134} ={}& -x_{177} - x_{179} x_{187} + x_{181}
\end{align*}
% END RESIDUAL 134
% BEGIN RESIDUAL 135
\begin{align*}
R_{135} ={}& 1 + x_{2} - x_{3} x_{18} - 4x_{176} x_{188} + x_{182}
\end{align*}
% END RESIDUAL 135
% BEGIN RESIDUAL 136
\begin{align*}
R_{136} ={}& 2 + x_{0}^{2} + 2x_{0} x_{2} - 2x_{0} x_{136} - 2x_{0} x_{143} \\
&{}  + 2x_{1} + x_{1}^{2} - 2x_{1} x_{170} - x_{2} + x_{2}^{2} \\
&{}  - 2x_{2} x_{143} + x_{4}^{2} - x_{4} x_{5} - 2x_{4} x_{161} - x_{6} x_{7} \\
&{}  + x_{8}^{2} - 2x_{8} x_{146} - 2x_{9} x_{10} - 2x_{9} x_{149} - 2x_{10} x_{152} \\
&{}  + x_{11}^{2} - 2x_{11} x_{155} + x_{12}^{2} - 2x_{12} x_{158} + x_{13}^{2} \\
&{}  - 2x_{13} x_{164} + x_{14}^{2} - 2x_{14} x_{167} + x_{15}^{2} - 2x_{15} x_{173} \\
&{}  - 4x_{16} x_{176} + x_{23} + x_{131}^{2} + x_{131} x_{136} - x_{131} x_{143} \\
&{}  - x_{131} x_{146} - x_{131} x_{149} - x_{131} x_{152} - x_{131} x_{155} - x_{131} x_{158} \\
&{}  + x_{131} x_{161} - x_{131} x_{164} - x_{131} x_{167} - x_{131} x_{173} - x_{131} x_{176} \\
&{}  - x_{131} x_{192} + x_{133}^{2} - x_{134}^{2} + x_{135} - x_{137} \\
&{}  + x_{137}^{2} - x_{142} - x_{144} + x_{144}^{2} - x_{145} \\
&{}  - x_{147} + x_{147}^{2} - x_{148} - x_{150} + x_{150}^{2} \\
&{}  - x_{151} - x_{153} + x_{153}^{2} - x_{154} - x_{156} \\
&{}  + x_{156}^{2} - x_{157} - x_{159} + x_{159}^{2} + x_{160} \\
&{}  - x_{162} + x_{162}^{2} - x_{163} - x_{165} + x_{165}^{2} \\
&{}  - x_{166} - x_{168} + x_{168}^{2} + x_{171}^{2} - x_{172} \\
&{}  - x_{174} + x_{174}^{2} - x_{175} + x_{176}^{2} - x_{177} \\
&{}  + x_{177}^{2} + x_{191} + x_{192} - x_{195}^{2} - x_{196}^{2} \\
&{}  - x_{197} - x_{197}^{2} + x_{210}
\end{align*}
% END RESIDUAL 136
% BEGIN RESIDUAL 137
\begin{align*}
R_{137} ={}& -x_{96} x_{193} + x_{120}
\end{align*}
% END RESIDUAL 137
% BEGIN RESIDUAL 138
\begin{align*}
R_{138} ={}& -x_{113} x_{194} + x_{120}
\end{align*}
% END RESIDUAL 138
% BEGIN RESIDUAL 139
\begin{align*}
R_{139} ={}& -1 - 4x_{16}^{2} - x_{136}^{2} - x_{143}^{2} - x_{146}^{2} \\
&{}  - x_{149}^{2} - x_{152}^{2} - x_{155}^{2} - x_{158}^{2} - x_{161}^{2} \\
&{}  - x_{164}^{2} - x_{167}^{2} - x_{170}^{2} - x_{173}^{2} + x_{176}^{2} \\
&{}  - x_{195}^{2} - x_{196}^{2} - x_{197} - x_{197}^{2}
\end{align*}
% END RESIDUAL 139
% BEGIN RESIDUAL 140
\begin{align*}
R_{140} ={}& -2 - x_{1} + x_{6} x_{7} - 4x_{98}^{2} - 4x_{102}^{2} \\
&{}  - 4x_{105}^{2} - 4x_{108}^{2} - 4x_{111}^{2} - 4x_{115}^{2} - 4x_{118}^{2} \\
&{}  + x_{120}^{2} - 4x_{122}^{2} - 4x_{126}^{2} - 4x_{129}^{2} - x_{198}^{2} \\
&{}  - x_{199}^{2} + 3x_{200} - x_{200}^{2}
\end{align*}
% END RESIDUAL 140
% BEGIN RESIDUAL 141
\begin{align*}
R_{141} ={}& -x_{2} x_{12} + x_{201}
\end{align*}
% END RESIDUAL 141
% BEGIN RESIDUAL 142
\begin{align*}
R_{142} ={}& x_{94} x_{203} - 2x_{95} x_{143} - 2x_{100} x_{146} - 2x_{103} x_{149} - 2x_{106} x_{152} \\
&{}  - 2x_{109} x_{155} - 2x_{112} x_{158} - 2x_{116} x_{164} - 2x_{119} x_{167} - 2x_{124} x_{173} \\
&{}  - 2x_{127} x_{176} - 2x_{131} x_{203} + 2x_{143} x_{144} + 2x_{146} x_{147} + 2x_{149} x_{150} \\
&{}  + 2x_{152} x_{153} + 2x_{155} x_{156} + 2x_{158} x_{159} + 2x_{164} x_{165} + 2x_{167} x_{168} \\
&{}  + 2x_{173} x_{174} + 2x_{176} x_{177} + x_{202}
\end{align*}
% END RESIDUAL 142
% BEGIN RESIDUAL 143
\begin{align*}
R_{143} ={}& -x_{94} x_{202} + x_{95}^{2} - 2x_{95} x_{142} + x_{100}^{2} - 2x_{100} x_{145} \\
&{}  + x_{103}^{2} - 2x_{103} x_{148} + x_{106}^{2} - 2x_{106} x_{151} + x_{109}^{2} \\
&{}  - 2x_{109} x_{154} + x_{112}^{2} - 2x_{112} x_{157} + x_{116}^{2} - 2x_{116} x_{163} \\
&{}  + x_{119}^{2} - 2x_{119} x_{166} + x_{124}^{2} - 2x_{124} x_{172} + x_{127}^{2} \\
&{}  - 2x_{127} x_{175} + x_{142}^{2} - x_{143}^{2} + x_{145}^{2} - x_{146}^{2} \\
&{}  + x_{148}^{2} - x_{149}^{2} + x_{151}^{2} - x_{152}^{2} + x_{154}^{2} \\
&{}  - x_{155}^{2} + x_{157}^{2} - x_{158}^{2} + x_{163}^{2} - x_{164}^{2} \\
&{}  + x_{166}^{2} - x_{167}^{2} + x_{172}^{2} - x_{173}^{2} + x_{175}^{2} \\
&{}  - x_{176}^{2} + x_{203}
\end{align*}
% END RESIDUAL 143
% BEGIN RESIDUAL 144
\begin{align*}
R_{144} ={}& x_{94} x_{205} - 2x_{96} x_{136} - 2x_{113} x_{161} - 2x_{120} x_{170} - 2x_{131} x_{205} \\
&{}  + 2x_{136} + 2x_{136} x_{137} + 2x_{161} + 2x_{161} x_{162} + 2x_{170} \\
&{}  + 2x_{170} x_{171} + x_{204} - x_{205}
\end{align*}
% END RESIDUAL 144
% BEGIN RESIDUAL 145
\begin{align*}
R_{145} ={}& -x_{94} x_{204} + x_{96}^{2} - 2x_{96} x_{135} + x_{113}^{2} - 2x_{113} x_{160} \\
&{}  + x_{120}^{2} - 2x_{120} x_{169} + x_{135}^{2} - x_{136}^{2} + x_{160}^{2} \\
&{}  - x_{161}^{2} + x_{169}^{2} - x_{170}^{2} + x_{204} + x_{205}
\end{align*}
% END RESIDUAL 145
% BEGIN RESIDUAL 146
\begin{align*}
R_{146} ={}& -1 - x_{2} - 2x_{2} x_{206} + x_{4}
\end{align*}
% END RESIDUAL 146
% BEGIN RESIDUAL 147
\begin{align*}
R_{147} ={}& -x_{2} x_{211} + x_{208}
\end{align*}
% END RESIDUAL 147
% BEGIN RESIDUAL 148
\begin{align*}
R_{148} ={}& -1 - 2x_{2} x_{208} + x_{207}^{2} - x_{208}^{2}
\end{align*}
% END RESIDUAL 148
% BEGIN RESIDUAL 149
\begin{align*}
R_{149} ={}& x_{2} - 2x_{209} x_{211} - x_{211}
\end{align*}
% END RESIDUAL 149
% BEGIN RESIDUAL 150
\begin{align*}
R_{150} ={}& -x_{0}^{2} + x_{210}
\end{align*}
% END RESIDUAL 150
% BEGIN RESIDUAL 151
\begin{align*}
R_{151} ={}& -x_{210}^{2} + x_{211}
\end{align*}
% END RESIDUAL 151
}


\end{document}
